\documentclass[11pt]{article}
\usepackage{iftex}
\ifPDFTeX
  \usepackage[T1]{fontenc}
  \usepackage{lmodern}
\fi
\usepackage{amsmath,amssymb,amsthm}
\usepackage[margin=1in]{geometry}
\usepackage{booktabs,longtable,array}
\usepackage[numbers,sort&compress]{natbib}
\usepackage{microtype}
\usepackage{xurl}
\usepackage{hyperref}
\hypersetup{colorlinks=true,linkcolor=blue,citecolor=blue,urlcolor=blue,
  pdftitle={The Möbius function at N - p},
  pdfauthor={Jizhou Guo}}
\Urlmuskip=0mu plus 1mu
\setlength{\emergencystretch}{2em}
\newtheorem{theorem}{Theorem}[section]
\newtheorem{lemma}[theorem]{Lemma}
\newtheorem{proposition}[theorem]{Proposition}
\newtheorem{corollary}[theorem]{Corollary}
\newtheorem{inputthm}[theorem]{Input}
\theoremstyle{definition}
\newtheorem{definition}[theorem]{Definition}
\newcommand{\e}{\mathrm e}
\newcommand{\E}{\mathbb E}
\newcommand{\Prob}{\mathbb P}
\newcommand{\Z}{\mathbb Z}
\newcommand{\N}{\mathbb N}
\newcommand{\Q}{\mathbb Q}
\newcommand{\C}{\mathbb C}
\newcommand{\R}{\mathbb R}
\newcommand{\ind}{\mathbf 1}
\newcommand{\Dpret}{\mathbb D}
\newcommand{\ph}{\varphi}
\newcommand{\leanname}[1]{\nolinkurl{#1}}
\newcommand{\norm}[1]{\left\lVert#1\right\rVert}
\newcommand{\abs}[1]{\left\lvert#1\right\rvert}
\DeclareMathOperator{\rad}{rad}
\DeclareMathOperator{\tr}{tr}
\DeclareMathOperator{\HS}{HS}
\DeclareMathOperator{\Rea}{Re}
\DeclareMathOperator{\cond}{cond}

\DeclareMathOperator{\supp}{supp}
\newcommand{\GN}{\mathcal N}
\newcommand{\T}{\mathbb T}
\newcommand{\PP}{\mathbb P}
\newcommand{\Aset}{\mathcal A}
\newcommand{\CC}{\mathcal C}
\newcommand{\QQ}{Q_Y^{\mathrm{all}}}
\newcommand{\Qmaj}{Q_Y^{\mathrm{maj},\GN}}
\newcommand{\tauN}{\tau_{\GN}}
\newcommand{\gN}{g_{\GN}}
\newcommand{\MN}{\mathfrak M_{\GN}}
\newcommand{\Mone}{\mathfrak M_1}
\newcommand{\LWW}[1]{\cite[#1]{LWW}}
\theoremstyle{remark}
\newtheorem{remark}[theorem]{Remark}
\title{The M\"obius function at $N-p$}
\author{Jizhou Guo\\
  \href{mailto:mitsuha2021b@gmail.com}{\texttt{mitsuha2021b@gmail.com}}}
\date{October 2026}
\numberwithin{equation}{section}
\begin{document}
\maketitle
\begin{abstract}
For even $\mathcal N$ tending to infinity, we deduce cancellation of
$\mu(\mathcal N-p)$ over primes $p<\mathcal N$, and equal densities of its
two nonzero signs, from a zero-free region implied by the $7/8$
half-plane of OpenAI and the arguments of L\"u, Wang and Wu as written.
We give the changes to their determinant method for this varying shift.
For each fixed $C$, Iwaniec's 1974 zero-free region and Burgess's bound
supply the analytic input when $\operatorname{rad}(\mathcal N)\le
(\log\mathcal N)^C$.
Assuming the additional arguments of L\"u, Wang and Wu as written, we
also obtain cancellation uniformly for $0<|h|\le(\log x)^C$ in
$\mu(p+h)$ on positive arguments. The corresponding Liouville statements
have sign densities $1/2+o(1)$.
\end{abstract}
\noindent\textit{2020 Mathematics Subject Classification.}
Primary 11N37; Secondary 11P32, 11N05.

\section{Introduction}\label{sec:intro}
All logarithms are natural. The M\"obius function is denoted by $\mu$,
and $\lambda(n)=(-1)^{\Omega(n)}$ is the Liouville function on positive
integers. Here $\Omega$ counts prime factors with multiplicity.
Write $\rad(n)=\prod_{p\mid n}p$ and let $\pi(x)$ count the primes at
most $x$. Put
\begin{equation}\label{eq:squarefree-product}
 A_{\GN}=\prod_{p\nmid\GN}\left(1-\frac1{p(p-1)}\right),\qquad
 A_h=\prod_{p\nmid h}\left(1-\frac1{p(p-1)}\right).
\end{equation}
These products are at least Artin's constant $0.3739\ldots$.

We specify the inputs to the statements. The assertion $\mathrm H2'$ is
that the arguments of LWW Sections~4.2--4.7, including the proofs of
Propositions~4.3--4.5 and Proposition~4.2, and Appendix~A,
(A.2)--(A.4), are correct as written
\LWW{pp.~17--39, 48--50}. The assertion $\mathrm H1$ is
\begin{equation}\label{eq:H1}
 L(s,\chi)\ne0\quad(\Re s>7/8)
 \quad\hbox{for every Dirichlet character $\chi$},
\end{equation}
with the principal pole at $s=1$ allowed. This is the Dirichlet
conclusion of \cite[Section~1, Theorem~1.1, p.~4]{OpenAIQ78}.
For $L=\log\GN$ and fixed $A_0$, write
$\mathrm H1^\star(C_{\mathrm{zf},1},C_{\mathrm{zf},2};A_1)$ at $\GN$
for the following assertion: for every $k\le L^{A_0}$ and every primitive
character $\chi^*$ of conductor $q^*\mid k\GN$ with $q^*>L^{A_1}$,
\begin{equation}\label{eq:Hstar}
 L(\sigma+it,\chi^*)\ne0\quad\left(
 \sigma\ge1-C_{\mathrm{zf},1}\frac{\log L}{L},\quad
 |t|\le L^{C_{\mathrm{zf},2}}\right).
\end{equation}
For every fixed choice of these constants, \eqref{eq:H1} implies
\eqref{eq:Hstar} at all sufficiently large $\GN$, since
$1-C_{\mathrm{zf},1}\log L/L>7/8$.

\begin{theorem}[Goldbach shift]\label{thm:G}
Assume $\mathrm H2'$ and $\mathrm H1^\star$ at every sufficiently large
even $\GN$, with the fixed constants chosen in
Theorem~\ref{thm:A} and Section~\ref{sec:assembly}.
Then, as $\GN\to\infty$ through even integers,
\begin{equation}\label{eq:G}
 \sum_{p<\GN}\mu(\GN-p)=o(\pi(\GN)),\qquad
 \#\{p<\GN:\mu(\GN-p)=\kappa\}
 =\left(\frac{A_{\GN}}2+o(1)\right)\pi(\GN)
 \quad(\kappa=\pm1).
\end{equation}
In particular these conclusions are deduced from $\mathrm H2'$ and
\eqref{eq:H1}. The errors are uniform in even $\GN$ and ineffective.
\end{theorem}

\begin{theorem}[Small radical]\label{thm:S}
Assume $\mathrm H2'$. For each fixed $C>0$, the conclusions
\eqref{eq:G} hold as $\GN\to\infty$ through even integers satisfying
$\rad(\GN)\le(\log\GN)^C$, using Iwaniec's region
\eqref{eq:Iwaniec} and the Burgess bound \eqref{eq:Burgess}.
The errors are uniform in this class for fixed $C$ and ineffective.
\end{theorem}

\begin{theorem}[Small shifts]\label{thm:H}
Assume the arguments of LWW as written: $\mathrm H2'$ and, in addition,
Sections~2.3, 3, 5 and 6 \LWW{pp.~6--14, 39--47}.
For each fixed $C>0$, uniformly in integers $0<|h|\le(\log x)^C$,
\begin{equation}\label{eq:H}
 \sum_{\substack{p\le x\\p+h\ge1}}\mu(p+h)=o(\pi(x)),\qquad
 \#\{p\le x:p+h\ge1,\ \mu(p+h)=\kappa\}
 =\left(\frac{A_h}2+o(1)\right)\pi(x)
 \quad(\kappa=\pm1).
\end{equation}
The errors are ineffective.
\end{theorem}

\begin{theorem}[Liouville versions]\label{thm:lambda}
Under the respective hypotheses of Theorems~\ref{thm:G}, \ref{thm:S}
and~\ref{thm:H}, their conclusions hold with $\lambda$ in place of
$\mu$ and each sign density replaced by $1/2+o(1)$.
The domain $p+h\ge1$ is retained in the small-shift statement.
\end{theorem}

\paragraph{Background.}\label{par:background}
In their preprint submitted on 14 October 2026, L\"u, Wang and Wu
\cite{LWW} state in Theorem~1 that
$\sum_{p\le x}\mu(p+2)=o(\pi(x))$, and in Theorem~2 give the sign
densities $A_2/2$ among all primes. Their abstract describes the
conjecture as ``dating back at least to Hildebrand's work in 1989'';
Section~1 attributes its appearance in print to Hildebrand's 1989
paper \LWW{p.~3}. Their Section~4.4 describes the construction of its
auxiliary operators as similar to Sections~4.1--4.3 of OpenAI's
\emph{The Poisson--Dirichlet law for prime predecessors}, dated
24 September 2026 \cite{OpenAIPD}; their Appendix~A says that its
root-counting proof follows the argument of that paper's Lemma~3.3,
with an explicit error for arbitrary squarefree $S$
\LWW{pp.~26, 47}.
The author's related papers, \emph{Shifted primes with a prescribed
parity of the number of prime factors} \cite{GuoSPP} and
\emph{Sums of two rough numbers with prescribed Liouville signs}
\cite{GuoRGP}, are dated 9 October 2026 in their repository entries.

\paragraph{What is used from LWW.}\label{par:lww}
We use Proposition~2.1 (the sifted Vaughan decomposition),
Proposition~2.2 with (2.5) and (2.11) (twisted sifted M\"obius sums),
Proposition~2.3 (weight removal), Proposition~3.1 with (3.1), (3.4),
(3.5) (Type~I and Bombieri--Vinogradov), Lemma~4.1 with (4.6), and
Proposition~6.1 with (6.3), as numbered inputs
\LWW{pp.~4--16, 46--47}. We rework the expansion (4.3)--(4.5), the
determinant equation (4.8), the block geometry, and the root-counting
modifications of Lemma~A.1 and (A.5). The arguments of
Sections~4.2--4.7 and Appendix~A are used as written after the changes
in Appendix~\ref{app:dictionary}.
In particular the proof of Proposition~4.2, including
Propositions~4.3--4.5, is the input $\mathrm H2'$.
The small-shift theorem also uses Lemmas~5.1--5.2 and the proof of
Proposition~5.3 \LWW{pp.~40--45}; its assumption therefore includes
more than $\mathrm H2'$. Formula~(5.15) \LWW{p.~44} is used only in
Remark~\ref{rem:other-contour}.

From \emph{The Quasi-Riemann Hypothesis}, dated 30 September 2026,
we use the zero-free half-plane \eqref{eq:H1}.
Its implication for \eqref{eq:Hstar} supplies an analytic input to
Theorem~\ref{thm:G}. Remark~\ref{rem:other-contour} records a
separate consequence of \eqref{eq:H1} at polynomial height. For the small-radical class we use
Iwaniec's Theorem~2 (1974, p.~97), stated in
Input~\ref{input:I}, and the $r=2$ Burgess bound for arbitrary
moduli in the form of Heath-Brown's introduction, equation~(1)
\cite{HeathBrown2012}. The real-character exception is removed in
Lemma~\ref{lem:Q}. The pointwise prime estimate needed here is
established in Lemma~\ref{lem:P}, rather than obtained from a
short-interval prime theorem.

\paragraph{Changes to LWW's argument.}
The determinant congruence decomposes the endpoint support into blocks
on the unit-coordinate chart of $\PP^1(\Z/\GN\Z)$.
In a basis of each index-$\GN$ lattice,
$\det(\gN z,\gN z')=\GN\det(z,z')$.
The matrix replacing $\operatorname{diag}(1,-2)$ is
$E_{\GN}=\operatorname{diag}(1,\GN)$, followed by an integral shear.
Lemma~\ref{lem:R} gives the root count with both coordinates units
modulo $\GN$ and with the normalized complementary ratio.
The major arcs reduce to the character-averaged energy
\eqref{eq:alpha-star} at polylogarithmic height.
In Type~I, $g=(r,\GN)$ is included in the extracted divisor before
applying the M\"obius identity, giving a reduced class modulo $2qr/g$.

\paragraph{Dependence and parameters.}
All uses of the LWW moment, memory, truncation and distinctness
arguments have exactly the scope specified above. The obstruction
in Proposition~\ref{prop:EC} explains a limitation of the classical
zero-free inputs for this method; the region \eqref{eq:Hstar} is a
sufficient input. We choose every precision and exponent before the
large-parameter limit. The coefficient-loss constant
$C_{\mathrm{coef}}$ and the zero-free constants
$C_{\mathrm{zf},1},C_{\mathrm{zf},2}$ are separate.
For weight removal the large parameter tends to infinity with a
fixed number $b$ of disjoint weights, and then $b$ tends to infinity.
Every $o(1)$ here is ineffective through Siegel--Walfisz or LWW
Proposition~2.2.

\paragraph{Outline.}
Sections~\ref{sec:setting}--\ref{sec:transport} give the determinant
reduction, blocks, root count and transport. Section~\ref{sec:major}
proves the major-arc energy estimate. Sections~\ref{sec:typeI}--\ref{sec:assembly}
complete the Goldbach statements. Sections~\ref{sec:small} and
\ref{sec:liouville} give the small-shift and Liouville statements.
Appendix~\ref{app:dictionary} records the changes to LWW by section,
equation and page.

\section{Setting and determinant reduction}\label{sec:setting}
\begin{definition}[Parameters and coefficients]\label{def:parameters}
For the Goldbach shift, $\GN$ is even and $L=\log\GN$.
Fix $B>2$, let $p\in(x,2x]$, and put
\begin{equation}\label{eq:ranges}
 \begin{gathered}
 \frac{\GN}{2L^B}\le x\le\frac{\GN}2,\quad
 Z=\exp(L^{0.24}),\quad U=V=x^{1/5},\\
 X=MN\asymp x,\quad M,N\ge x^{1/5}/2,\quad MN\le2x\le\GN.
 \end{gathered}
\end{equation}
The letter $N$ always denotes the length of the Vaughan variable $n$.
Write $R_Z(n)$ for the indicator that every prime divisor of $n$
exceeds $Z$. Fix $0<\rho<1$ and
\begin{equation}\label{eq:bands}
 \frac1{10}<a_1<\cdots<a_K<\frac15,\quad
 \mathcal P_i=\{p:\exp(L^{a_i})<p\le\exp(2L^{a_i}),\ p\nmid\GN\},
 \quad V_i=\sum_{p\in\mathcal P_i}\frac1p.
\end{equation}
Put $\Aset=\{p_1\cdots p_K:p_i\in\mathcal P_i\}$,
$\omega_i(n)=\#\{p\in\mathcal P_i:p\mid n\}$ and
$\omega(n)=\sum_i\omega_i(n)$. The weight of LWW (2.13), p.~7, is
\begin{equation}\label{eq:weight}
 W(n)=\prod_{i=1}^K\frac{\rho^{\omega_i(n)-1}\omega_i(n)}{V_i}.
\end{equation}
Let $J=\lfloor L^{0.01}\rfloor$, $R=\lceil L^{1/2}/2\rceil$,
and let $d_0,D$ have the meanings of LWW Section~4.2, pp.~17--21.
The real smooth dyadic cutoff $\eta$ is supported on $[1,4]$
(LWW Section~4.1, p.~15), so its product scale $Y$ has
$a,b\in[Y,4Y]\subset[Y,16Y]$. Define
\begin{equation}\label{eq:rectangle}
 \Omega=[d_0YM,16d_0YM]\times[N,2N].
\end{equation}
For $|v|\le L^{2B}$, the coefficients are
\begin{equation}\label{eq:coefficients}
 \begin{split}
 \alpha_m&=\mu(m)R_Z(m)\ind_{(m,\GN)=1}m^{iv},\\
 \beta_n&=c_V(n)R_Z(n)\ind_{(n,\GN)=1}n^{iv},\qquad
 c_V(n)=\sum_{\substack{\ell\mid n\\\ell>V}}\Lambda(\ell).
 \end{split}
\end{equation}
Thus $|\alpha_m|\le1$, $|\beta_n|\le L$, and the untwisted
coefficient $c_V(n)R_Z(n)\ind_{(n,\GN)=1}$ is nonnegative.
For a character $\xi$ of modulus dividing $k\GN$, put
\begin{equation}\label{eq:polynomials}
 \begin{split}
 B_\xi(t)&=\frac1{Y\prod_i V_i}
    \sum_{b\in\Aset}\eta(b/Y)e(\lambda_0b/Y)\xi(b)b^{it},\\
 M_\xi(t)&=\frac1M\sum_{M<m\le2M}\alpha_m\xi(m)m^{it},\qquad
 N_\xi(t)=\frac1N\sum_{N<n\le2N}\beta_n\xi(n)n^{it}.
 \end{split}
\end{equation}
Here $e(t)=\exp(2\pi it)$, $k\le L^{A_0}$,
$|\lambda_0|\le2L^{A_0}$, $T_0=\lceil L^{B_1}\rceil$, and
$|B_\xi(t)|\le C_B=16\|\eta\|_\infty/\prod_i V_i$.
We write the arc displacement as $\lambda_0$ to distinguish it from
the Liouville function.
The major-arc target, at a fixed requested precision $A$, is
\begin{equation}\tag{$\alpha^*$}\label{eq:alpha-star}
 \frac1{\ph(\GN)}\sum_{\chi\bmod\GN}
 \int_{|t|\le T_0}|B_{\chi\psi_2}(t)M_{\chi\psi}(t)N_{\chi\psi}(t)|^2\,dt
 \ll_A\frac{L^{-A}}{\GN}.
\end{equation}
\stepcounter{equation}
The fixed twists $\psi,\psi_2$ are those constructed in
Section~\ref{subsec:arc-reduction}. The trivial bound is
$L^{O(1)}/\GN$; $1/\GN$ is the normalization.
\end{definition}

\begin{remark}[Ranges]\label{rem:ranges}
Primes $p\le\GN/L^B$ contribute $O(\GN/L^B)$ to the bounded target
sum, and $O(\GN L^{1-B})$ with $\log p$. The rest splits into
$O(\log L)$ dyadic intervals in \eqref{eq:ranges}.
Before separating the product cutoff,
$\GN-mn\in[\GN-2x,\GN-x)\cap(0,\GN)$.
For fixed $a\asymp Y$ the positive cofactor belongs to
\begin{equation}\label{eq:fixed-cofactor}
 \left[\frac{\GN-2x}{a},\frac{\GN-x}{a}\right],
\end{equation}
an interval of length $x/a\asymp X/Y$. As $a$ varies, its centre
varies, and the common positive support has $O(\GN/Y)$ integers.
Terms with $0<\GN-mn<Y$ vanish, since $W$ requires a factor from
every band. No dyadic condition on $\GN-p$ is needed.
\end{remark}

\begin{proposition}[Vaughan reduction]\label{prop:vaughan}
Define $F(n)=\mu(\GN-n)W(\GN-n)$ for $1\le n<\GN$ and
$F(n)=0$ otherwise. LWW Proposition~2.1 \LWW{p.~4} applies to this
bounded function. Its Type~II sum is
\begin{equation}\label{eq:T3}
 T_3(F)=\sum_{\substack{m>U,n>V\\x<mn\le2x}}
 \mu(m)R_Z(m)c_V(n)R_Z(n)\mu(\GN-mn)W(\GN-mn).
\end{equation}
The smoothing and Mellin separation in LWW Section~4.1, p.~14,
reduce it to rectangular sums
\begin{equation}\label{eq:Trect}
 T=\sum_{m,n}\alpha_m\beta_n F(mn),\qquad
 M,N\ge x^{1/5}/2,\quad MN\asymp x,\quad |v|\le L^{2B},
\end{equation}
with the coprimality restrictions justified below. The argument of
$\mu$ is positive wherever $F$ is nonzero.
\end{proposition}
\begin{proof}
Apply the stated decomposition and separate the cutoff in $mn/x$.
It concerns the prime variable, so this separation has the same
Mellin norm and coefficient bounds for $F$ as in LWW. On parts of a
rectangle where $mn\ge\GN$, the definition $F(mn)=0$ applies.
\end{proof}

\begin{lemma}[Puncturing and unit support]\label{lem:puncturing}
The punctured bands satisfy $V_i=\log2+o(1)$, uniformly in $\GN$.
The restrictions $(m,\GN)=(n,\GN)=1$ in \eqref{eq:coefficients}
cost $O_A(XL^{-A})$ for every fixed $A$.
The remaining variables are units modulo $k\GN$ whenever $k\le L^{A_0}$.
\end{lemma}
\begin{proof}
The removed reciprocal mass is at most
\begin{equation}\label{eq:puncture-mass}
 \sum_{\substack{p\mid\GN\\p>\exp(L^{0.1})}}\frac1p
 \le\omega_{\rm all}(\GN)e^{-L^{0.1}}\le Le^{-L^{0.1}},
\end{equation}
where $\omega_{\rm all}$ counts all distinct prime divisors.
Thus LWW (2.14)--(2.15), p.~7, apply with the same limits.
A common prime of a sifted variable and $\GN$ exceeds $Z$.
The number of discarded pairs is bounded, up to an absolute constant, by
\begin{equation}\label{eq:unit-discard}
 \sum_{\substack{p\mid\GN\\p>Z}}\frac Xp
 \le\frac{X\omega_{\rm all}(\GN)}Z\le XLe^{-L^{0.24}}.
\end{equation}
The bound $|\alpha_m\beta_n|\ll L$ absorbs this into the claimed
error. The restrictions preserve the coefficient and sifting
hypotheses of Proposition~4.2. Every prime factor of $k$ is below
$Z$, so the sifted variables are also units modulo $k$.
Every list prime and every $a,b\in\Aset$ is a unit modulo $\GN$ by
puncturing, and modulo $k$ because the bands lie above $L^{A_0}$.
\end{proof}

\begin{proposition}[The cofactor square]\label{prop:cofactor-square}
After a smooth dyadic partition $\eta(a/Y)$, the expansion of
\eqref{eq:Trect} has the form
\begin{equation}\label{eq:expansion-mu}
 T=\frac{(-1)^K}{\prod_i V_i}
 \sum_{h\ge1}\mu(h)\rho^{\omega(h)}
 \sum_{\substack{p_i\in\mathcal P_i,m,n\\mn+ah=\GN}}
 \alpha_m\beta_n+O(XLe^{-L^{0.1}}).
\end{equation}
Its dyadic pieces satisfy
\begin{equation}\label{eq:cauchy-all}
 |T_Y|^2\ll\frac{\GN}{Y}Q_Y^+\le\frac{\GN}{Y}\QQ,
 \qquad
 \QQ=\sum_{\substack{h\in\Z\\h\ne0}}
 \left|\frac1{\prod_i V_i}
 \sum_{\substack{a\in\Aset,m,n\\\GN-mn=ah}}
 \eta(a/Y)\alpha_m\beta_n\right|^2\ge0.
\end{equation}
Here $Q_Y^+$ is the same square restricted to $h\ge1$.
All cofactors in $\QQ$ satisfy $|h|\ll\GN/Y$.
\end{proposition}
\begin{proof}
If $\GN-mn=ah$ and $(a,h)=1$, then
$\mu(ah)=(-1)^K\mu(h)$ and $\omega(ah)=K+\omega(h)$.
Also $(a,m)=1$: a common list prime would divide $\GN$.
When $(a,h)>1$, some list prime has square dividing $\GN-mn$.
For fixed $m$ this is one residue class of $n$ modulo that square,
and the count in LWW (4.3), p.~15, gives $O(XLe^{-L^{0.1}})$.
The coprimality restriction can be removed in the expansion at the
same cost; its new terms have that square-divisibility property.
Cauchy--Schwarz uses $|\mu(h)\rho^{\omega(h)}|\le1$ and the
common range $O(\GN/Y)$ from Remark~\ref{rem:ranges}.
Each square is nonnegative, proving $Q_Y^+\le\QQ$.
On a supported rectangle $mn\ll MN\le\GN$, whence
$|\GN-mn|\ll\GN$. The cofactor $0$ is impossible because
$mn=\GN$ contradicts $(mn,\GN)=1$ for $\GN>1$.
Only this upper-bound quadratic form introduces negative cofactors;
$F$ and the expansion of its target always use positive arguments.
\end{proof}

\begin{lemma}[Disjoint lists and determinant equation]\label{lem:determinant}
Terms with $(a,b)>1$ may be removed from $\QQ$ at a cost
$O_A(\GN YL^{-A})$. For $(a,b)=1$ the exact determinant equation is
\begin{equation}\label{eq:determinant}
 b(\GN-mn)=a(\GN-rs)
 \quad\Longleftrightarrow\quad
 t:=bmn-ars=\GN(b-a).
\end{equation}
Equality on the left is equality of the two integer cofactors,
with no positivity condition. In particular $\GN\mid t$ and
$|t|\ll\GN Y$, while $t'=t/\GN=b-a$ has size $O(Y)$.
\end{lemma}
\begin{proof}
The divisor-square estimate used in LWW Lemma~4.1 and (4.6), p.~16,
has the form
\begin{equation}\label{eq:divisor-square}
 \sum_{h\in I_\ell}\tau(\GN-\ell h)^2\ll(X/Y)L^3,\qquad |I_\ell|\ll X/Y
\end{equation}
where $I_\ell$ is the translated cofactor interval for fixed
$\ell\asymp Y$, with $\GN-\ell h$ in the positive dyadic
product range. The argument depends on
$d\mid\GN-\ell h\Rightarrow(d,\ell)\mid\GN$; since
$(\ell,\GN)=1$, it gives $(d,\ell)=1$, exactly the needed hypothesis.
For the all-integer square use the full interval $|h|\ll\GN/Y$
and $|\GN-\ell h|$ in its absolute counts. Its length costs at most
$\GN/X\ll L^B$ relative to the displayed local estimate.
LWW's divisor argument and the common-list saving then give the
stated $\GN YL^{-A}$ after increasing the fixed precision.
Reducing \eqref{eq:determinant} modulo $a$ and $b$, and using
$(a,b)=1$, proves divisibility by $a,b$ of the corresponding
numerators and equality of their integer quotients.
\end{proof}

\section{Blocks and relative coordinates}\label{sec:blocks}
Endpoint positions, with the list data of LWW attached, are
\begin{equation}\label{eq:positions}
 P=(Dbm,s)^T,\qquad Q=(Dar,n)^T,\qquad
 \det(P,Q)=D(bmn-ars)=Dt.
\end{equation}
All endpoint coordinates, and $D,a,b$, are units modulo $\GN$.

\subsection{The major-arc congruence}\label{subsec:kernel}
Let $\Mone$ be LWW's union of arcs with denominators
$k\le L^{A_0}$ and width $2L^{A_0}/Y$
\LWW{Section~4.1, p.~16}. Define
\begin{equation}\label{eq:arc-kernel}
 \begin{split}
 \MN&=\{\theta\in\T:\GN\theta\in\Mone\},\\
 H_{\GN}(t;a,b)&=\psi_1(t/(\GN Y))
 \int_{\MN}e(\theta(t-\GN(b-a)))\,d\theta,\\
 \Qmaj&=\prod_iV_i^{-2}
 \sum_{a,b,m,n,r,s}\eta(a/Y)\eta(b/Y)
 \alpha_m\overline{\alpha_r}\beta_n\overline{\beta_s}
 H_{\GN}(bmn-ars;a,b).
 \end{split}
\end{equation}
The cutoff $\psi_1$ is the fixed smooth cutoff of LWW, equal to $1$
on $[-4,4]$.
\begin{lemma}[Covering map]\label{lem:cover}
For every integer $t$,
\begin{equation}\label{eq:cover}
 \int_{\MN}e(\theta(t-\GN(b-a)))\,d\theta
 =\ind_{\GN\mid t}\int_{\Mone}e(\phi(t/\GN-b+a))\,d\phi.
\end{equation}
In particular $|\MN|=|\Mone|\ll L^{3A_0}/Y$.
\end{lemma}
\begin{proof}
For a fundamental domain of $\Mone$, the change of variables gives
\[
 \int_{\MN}F(\theta)\,d\theta
 =\frac1{\GN}\sum_{j\bmod\GN}\int_{\Mone}F((\phi+j)/\GN)\,d\phi.
\]
For the integrand in \eqref{eq:cover},
\[
 F((\phi+j)/\GN)
 =e(\phi t/\GN)e(jt/\GN)e(-\phi(b-a))e(-j(b-a)).
\]
The last factor is $1$ and the $j$-sum is
$\GN\ind_{\GN\mid t}$. This also makes the identity independent of
the chosen representatives.
\end{proof}
The raw-minus-major kernel is
\begin{equation}\label{eq:K}
 \begin{split}
 K_{\GN}(P,Q)={}&\frac{d_0}{D}\eta(a/Y)\eta(b/Y)\psi_1(t/(\GN Y))\\
 &\cdot\left\{\ind_{t=\GN(b-a)}
 -\int_{\MN}e(\theta(t-\GN(b-a)))\,d\theta\right\}
 \rho^{(\omega(P_1)+\omega(Q_1)-2K(J+1))/2}.
 \end{split}
\end{equation}
Both terms vanish unless $\GN\mid t$. On that congruence the brace is
\begin{equation}\label{eq:reduced-brace}
 \ind_{t'=b-a}-\int_{\Mone}e(\phi(t'-b+a))\,d\phi,
 \qquad t'=t/\GN,
\end{equation}
with cutoff $\psi_1(t'/Y)$, the transformed kernel of LWW
Section~4.2, p.~18.

\subsection{Lines and their sizes}\label{subsec:lines}
\begin{definition}[Unit endpoints and labelled blocks]\label{def:blocks}
Let $\CC_{\GN}$ consist of $P\in\Omega$ primitive in $\Z^2$
with $(P_1P_2,\GN)=1$. For a unit $\varrho\bmod\GN$, put
\begin{equation}\label{eq:lattice}
 \Lambda_\varrho=\{(u,v)\in\Z^2:v\equiv\varrho u\pmod{\GN}\}.
\end{equation}
This has basis $(1,\varrho),(0,\GN)$ and index $\GN$.
The slope class of a unit-coordinate endpoint is
$\varrho(P)=P_2\overline{P_1}\bmod\GN$.
\end{definition}
\begin{lemma}[Block diagonality]\label{lem:blocks}
For unit-coordinate endpoints,
\begin{equation}\label{eq:disjoint-union}
 \GN\mid\det(P,Q)\ \Longleftrightarrow\ \varrho(P)=\varrho(Q),
 \qquad
 \CC_{\GN}=\coprod_{\varrho\in(\Z/\GN\Z)^\times}
      (\CC_{\GN}\cap\Lambda_\varrho).
\end{equation}
Thus the endpoint kernel is block diagonal. Each full lattice is
regarded as a labelled copy when the state space is enlarged.
\end{lemma}
\begin{proof}
Invert $P_1Q_1$ in the determinant congruence. Since $(D,\GN)=1$,
$\GN\mid t$ is equivalent to $\GN\mid\det(P,Q)$.
The unit restriction makes the slope class unique.
\end{proof}
\begin{lemma}[Block size]\label{lem:block-size}
The number of vectors in $\Omega\cap\Lambda_\varrho$ primitive in
$\Lambda_\varrho$ is
\begin{equation}\label{eq:block-size}
 O\left(1+\frac{d_0YMN}{\GN}\right).
\end{equation}
In the application $d_0YMN/\GN\to\infty$ and $MN\le\GN$, so this is
$O(d_0YMN/\GN)\le O(d_0Y)$, uniformly in $\varrho$.
\end{lemma}
\begin{proof}
If the block is occupied, choose a lattice-primitive $P_0$ in it.
For any $Q$ in the block,
\[
 \det(P_0,Q)\in\GN\Z,\qquad
 |\det(P_0,Q)|\le2\cdot16d_0YM\cdot2N
 =64d_0YMN\le64d_0Y\GN.
\]
There are $O(1+d_0YMN/\GN)$ possible determinants. For each, two
solutions differ by an element of
$\Q P_0\cap\Lambda_\varrho=\Z P_0$; their first coordinates permit
at most $16$ solutions. Here $\Q$ denotes the rational numbers.
In the application $\GN/(d_0YMN)\ll L^B/(d_0Y)\to0$.
\end{proof}
This partition supplies pieces of the size of LWW's slope intervals
\LWW{Section~4.6, p.~36}. Physical slope locality is replaced by
relative-coordinate bounds.

\subsection{A normalized basis}\label{subsec:basis}
Let $P=(u,v)^T\in\CC_{\GN}$, and choose $ud-vc=1$, $0\le c<u$.
Define
\begin{equation}\label{eq:basis}
 \begin{split}
 c'&\equiv\GN c\pmod u,\quad0\le c'<u,
 \qquad c'v\equiv-\GN\pmod u,\quad \beta_P=c'/u,\\
 d''&=(\GN+c'v)/u,\qquad P^*=(c',d'')^T,\qquad
 \gN(P)=\begin{pmatrix}u&c'\\v&d''\end{pmatrix},\quad\det\gN(P)=\GN.
 \end{split}
\end{equation}
Equivalently
\begin{equation}\label{eq:shear}
 P^*=\GN(c,d)^T-\lfloor\GN c/u\rfloor P,\qquad
 \gN(P)=gE_{\GN}U_0,\quad
 g\in\operatorname{SL}_2(\Z),\quad E_{\GN}=\operatorname{diag}(1,\GN),
\end{equation}
where $U_0$ is the corresponding unimodular shear.
\begin{lemma}[Relative coordinates]\label{lem:coordinates}
Writing $\varrho=\varrho(P)$ and $Q=\gN(P)z$, one has:
\begin{enumerate}
\item $\Lambda_\varrho=\Z P\oplus\Z P^*$.
\item With $\tauN(z)=z_1+\beta_Pz_2$,
\begin{equation}\label{eq:coordinates}
 \det(P,Q)=\GN z_2,\quad Q_1=u\tauN(z),\quad
 Q_2=v\tauN(z)+\frac{\GN z_2}{u}.
\end{equation}
\item $Q$ is primitive in $\Lambda_\varrho$ precisely when
$(z_1,z_2)=1$. It is primitive in $\Z^2$ precisely when also
$(Q_1,\GN)=1$.
\item For primitive $z$ and every prime $p\nmid\GN$,
\begin{equation}\label{eq:kernel-line}
 p\mid Q_1\quad\Longleftrightarrow\quad
 [z]_p=\ker(u,c')\in\PP^1(\mathbb F_p).
\end{equation}
\item For $Q\in\Omega$,
\begin{equation}\label{eq:z-bound}
 |z_2|\le\frac{64d_0YMN}{\GN}\le64d_0Y,\qquad
 |z_1|\le16+|z_2|\le2^7d_0Y.
\end{equation}
Conversely every integer $z$ gives $Q\in\Lambda_\varrho$.
\end{enumerate}
\end{lemma}
\begin{proof}
The two columns belong to $\Lambda_\varrho$ and have its determinant.
Formula~\eqref{eq:coordinates} follows from
$d''=\GN/u+\beta_Pv$.
The lattice bijection proves lattice-primitivity.
At $p\mid\GN$ the matrix has rank one, with unit-coordinate column
$P$, and $Q\equiv0\pmod p$ precisely when $p\mid Q_1$.
At $p\nmid\GN$ the matrix is invertible, so $Q\equiv0$ precisely
when $z\equiv0$. Also $(u,c')=1$, since
$c'v\equiv-\GN\pmod u$ and $(u,\GN)=1$.
This proves the kernel assertion for primitive $z$.
Finally use the determinant bound in Lemma~\ref{lem:block-size},
$Q_1/u\le16$ and $0\le\beta_P<1$.
\end{proof}
The continuous extension and its comparison with LWW are
\begin{equation}\label{eq:continuous}
 \gN(u,v,\beta)=\begin{pmatrix}u&\beta u\\v&\GN/u+\beta v\end{pmatrix},
 \quad\tauN(z)=z_1+\beta z_2,
 \quad
 g'(u,v,r)=\begin{pmatrix}u&-2ur\\(u-v)/2&r(v-u)+1/u\end{pmatrix}.
\end{equation}
Thus $-2r$ is replaced by $\beta$ and $1/u$ by $\GN/u$; both
entries are retained exactly. The ratio in LWW (4.26), p.~21,
becomes $\beta_P=\{\GN r_P\}$. Multiplication by $\GN$ preserves
Lebesgue measure on $\T$.

\subsection{Enlarged states and intrinsic ratios}\label{subsec:enlargement}
\begin{definition}[States and good-state ratios]\label{def:good-ratio}
The ambient position space is the disjoint union of labelled sets
\begin{equation}\label{eq:states}
 \mathcal S'_\varrho=
 \{(\varrho,Q):Q\in\Omega\cap\Lambda_\varrho,
                  \ Q\text{ primitive in }\Lambda_\varrho\},
\end{equation}
with LWW's lists attached. For such a state choose $W\in\Lambda_\varrho$
with $\det(Q,W)=\GN$ and put
\begin{equation}\label{eq:intrinsic-ratio}
 \beta_Q=\frac{W_1}{Q_1}\pmod1.
\end{equation}
The good-state tests of LWW (4.14)--(4.16), pp.~18--19, use this
ratio at every state.
\end{definition}
\begin{lemma}[Intrinsic ratio and endpoint compression]\label{lem:compression}
Definition~\ref{def:good-ratio} is independent of the complement and
of the chosen block basis. In relative coordinates,
\begin{equation}\label{eq:ratio-relative}
 \beta_Q\equiv\frac{\tauN(w)}{\tauN(z)}\pmod1
 \quad\bigl(Q=\gN(P)z, W=\gN(P)w, \det(z,w)=1\bigr).
\end{equation}
For a unit endpoint it equals $\{\GN r_Q\}$.
LWW's test that $z_1$ is odd is replaced by endpoint support
$(Q_1Q_2,\GN)=1$, which is not a parity statement about the basis used here.
The restricted-root trace in LWW Section~4.6 is valid in this space.
\end{lemma}
\begin{proof}
A lattice-primitive $Q$ has a determinant-$\GN$ complement; any two
such complements differ by an integer multiple of $Q$.
Two oriented bases differ by $\operatorname{SL}_2(\Z)$, proving
\eqref{eq:ratio-relative} and basis independence.
At a unit endpoint the normalized complement in \eqref{eq:basis}
gives $\{\GN r_Q\}$. A nonunit vector can belong to several full
lattices, giving a different labelled state in each. Such a vector
is not primitive in $\Z^2$, by Lemma~\ref{lem:coordinates}; the
endpoint vector vanishes there because its coefficients and lists
are supported on units. For example, $\GN=2$, $P=(3,1)$,
$P^*=(1,1)$ and $z=(0,1)$ give a unit endpoint with $z_1$ even.
Thus no assertion $z_1$ odd is made after this shear.

What the moment argument uses is a completion at a root, the
primitive-coordinate equivalence
\[
 D\mid\det(z,z')\ \Longleftrightarrow\ [z]_p=[z']_p
 \quad(p\mid D),
\]
for squarefree $D$, and the ratio at the current state.
All hold here. The diagonal operators $S,G$ preserve the endpoint
pairing under enlargement. For positive $B_0$, the matrix
\begin{equation}\label{eq:compressed-gram}
 \bigl(\langle u_P,B_0^R u_Q\rangle\bigr)
       _{P,Q\in\CC_{\GN}\cap\Lambda_\varrho}
 \succeq0,\qquad
 \lambda_{\max}\le d_\varrho:
 =\sum_{P\in\CC_{\GN}\cap\Lambda_\varrho}
        \langle u_P,B_0^R u_P\rangle
\end{equation}
is a principal submatrix of a Gram matrix.
Spectral H\"older in LWW (4.65), p.~36, is applied in the ambient
space before this restriction. The sum $\sum_\varrho d_\varrho$
is exactly the root average over $\Omega\cap\CC_{\GN}$ counted below.
No estimate for the full trace on added nonunit states is required.
The second good-state test in LWW (4.47)--(4.48), p.~30, is applied
at the source position $Q$, using $\beta_Q$.
For the exceptional-measure bound, multiplication by $\GN$ preserves
Haar measure; for the subsequent cube count we count $\beta$ directly.
\end{proof}

\section{Root counting}\label{sec:roots}
\begin{lemma}[Lemma R]\label{lem:R}
Put $U_0=d_0YM$, $V_0=N$, and fix $\delta>0$.
Suppose $U_0,V_0\ge x^\delta$, $\GN<U_0V_0$, and
$\log x=L+O_B(\log L)$. Let $S$ be odd and squarefree,
$(S,\GN)=1$, $S\le\exp(C_0L^{0.98})$, and let
$g_0\bmod S$ have determinant $\GN\bmod S$.
For intervals $I_1\subset[1,16]$, $I_2\subset[1,2]$ and
$I_3\subset[0,1]$,
\begin{equation}\label{eq:R-matrix}
 \begin{split}
 &\#\{P=(u,v)^T:(u,v)=1,\ (uv,\GN)=1,
       \,u/U_0\in I_1,\ v/V_0\in I_2,
       \,\beta_P\in I_3,\ \gN(P)\equiv g_0\pmod S\}\\
 &\hspace{5mm}=
 \frac{c_{\GN}U_0V_0|I_1||I_2||I_3|}
      {|\operatorname{SL}_2(\Z/S\Z)|}
       +O(U_0V_0x^{-\delta/800}),
 \end{split}
\end{equation}
where
\begin{equation}\label{eq:cN}
 c_{\GN}=\frac1{\zeta(2)}\prod_{p\mid\GN}\frac{p-1}{p+1}
 =\frac1{\zeta(2)}\prod_{p\mid\GN}
       \frac{(1-1/p)^2}{1-1/p^2},\qquad
 (\log\log\GN)^{-2}\ll c_{\GN}\le1.
\end{equation}
For prescribed lines $\ell_p\in\PP^1(\mathbb F_p)$ at $p\mid S$,
\begin{equation}\label{eq:R-lines}
 \begin{split}
 &(U_0V_0)^{-1}
 \#\{P:(u,v)=1,(uv,\GN)=1,\ u/U_0\in I_1,\ v/V_0\in I_2,\\
 &\hspace{35mm}\beta_P\in I_3,\ \ker(\gN(P)_{1,*}\bmod p)=\ell_p\ (p\mid S)\}\\
 &\hspace{8mm}=c_{\GN}|I_1||I_2||I_3|
          \prod_{p\mid S}\frac1{p+1}+O(x^{-c\delta})
 \end{split}
\end{equation}
for an absolute $c>0$, possibly smaller than $1/800$.
All errors are uniform in the intervals, matrix class and $\GN$
in the stated dyadic regime.
\end{lemma}
\begin{proof}
We use LWW Appendix~A, (A.2)--(A.4), pp.~48--50, and give the
changes to Lemma~A.1 and (A.5), pp.~47--51.

\emph{The matrix and the shorter-coordinate case.}
The matrix \eqref{eq:basis} has determinant $\GN$ and $0\le c'<u$.
Suppose $U_0\le V_0$, and fix $u\equiv u_0\pmod S$ with
$(u,\GN)=1$. Replace LWW's set by
\begin{equation}\label{eq:Ru}
 \mathcal R_u=\{(c',v)\bmod uS:
       c'\equiv c'_0\ (S),\ v\equiv v_0\ (S),\
       c'v\equiv ud''_0-\GN\ (uS)\}.
\end{equation}
The right side in the last congruence is a unit at every prime
dividing $u$. At primes dividing both $u$ and $S$ this follows also
from $u_0d''_0-c'_0v_0\equiv\GN\not\equiv0\pmod p$.
Consequently the same count as in LWW gives
\begin{equation}\label{eq:Ru-size}
 \#\mathcal R_u=u\prod_{\substack{p\mid u\\p\nmid S}}(1-1/p),
 \qquad (u,c')=(u,v)=1.
\end{equation}

\emph{Exponential sums.}
Let $S_1=(u,S)$, $S_2=S/S_1$, and
$\gamma S_2\equiv1\pmod{uS_1}$. The CRT calculation of
LWW (A.3), p.~49, gives
\begin{equation}\label{eq:CRT}
 \left|\sum_{(c',v)\in\mathcal R_u}
        e\left(\frac{hc'+kv}{uS}\right)\right|
 \le\frac1{S_1}\sum_{j\bmod S_1}
   |K_{uS_1}(\gamma h+ju,\gamma k(ud''_0-\GN))|.
\end{equation}
Here $K_m(h,k)=\sum_{y\bmod m,(y,m)=1}e((hy+k\bar y)/m)$.
Since $S_1\mid u$ and $(u,\GN)=1$, $ud''_0-\GN$ is a unit
modulo $uS_1$. Thus
\[
 (\gamma h+ju,\gamma k(ud''_0-\GN),uS_1)\le S_1(h,k,u).
\]
The all-modulus fourth-moment bound (A.2), followed by (A.3) and
Fourier completion (A.4), gives rectangle discrepancy
$O(u^{0.99}S^2)$, independent of the rectangle.
No prime-modulus Kloosterman estimate is used.

\emph{The sieve on $v$.}
Expand
$\ind_{(v,\GN)=1}=\sum_{d\mid\rad(\GN)}\mu(d)\ind_{d\mid v}$.
There are at most $2^{\omega_{\rm all}(\GN)}=x^{o(1)}$ terms.
Every $d$ is coprime to $uS$; replacing $v$ by $dv_2$ replaces
the right side of \eqref{eq:Ru} by the unit
$\bar d(ud''_0-\GN)$ and the prescribed residue by
$v_2\equiv\bar d v_0\pmod S$. Therefore the same rectangle error
applies. Splitting into periods as in LWW Appendix~A, p.~50, the
count for this $d$ is
\begin{equation}\label{eq:v-sieve}
 \frac{V_0|I_2||I_3|}{dS^2}
 \prod_{\substack{p\mid u\\p\nmid S}}(1-1/p)
 +O\left(\left(\frac{V_0}{uSd}+1\right)u^{0.99}S^2\right).
\end{equation}
The sum with $\mu(d)$ multiplies the main term by
$\prod_{p\mid\GN}(1-1/p)$ and leaves error
$O(x^{o(1)}(V_0/(uS)+1)u^{0.99}S^2)$.

\emph{The sieve on $u$ and the local constant.}
Use
\[
 \prod_{\substack{p\mid u\\p\nmid S}}(1-1/p)
 =\sum_{\substack{\ell\mid u\\(\ell,S)=1}}\frac{\mu(\ell)}\ell,
 \qquad
 \ind_{(u,\GN)=1}=\sum_{d\mid\rad(\GN)}\mu(d)\ind_{d\mid u}.
\]
For $(\ell d,S)=1$, counting $u\in U_0I_1$ in its fixed residue
class gives
$U_0|I_1|/(\operatorname{lcm}(\ell,d)S)+O(1)$.
The double divisor series has local factors
\begin{equation}\label{eq:local-factors}
 \begin{cases}
 1-p^{-2},&p\nmid S\GN,\\
 (1-p^{-2})-p^{-1}(1-p^{-1})=1-p^{-1},&p\mid\GN,\\
 1,&p\mid S.
 \end{cases}
\end{equation}
To justify its tail without a factor $2^{\omega_{\rm all}(\GN)}$,
sum over $d$ before extending $\ell$:
\begin{equation}\label{eq:divisor-tail}
 \sum_{d\mid\rad(\GN)}\frac{\mu(d)}{\operatorname{lcm}(\ell,d)}
 =\begin{cases}
 0,&(\ell,\GN)>1,\\
 \ell^{-1}\prod_{p\mid\GN}(1-1/p),&(\ell,\GN)=1.
 \end{cases}
\end{equation}
Indeed a prime common to $\ell$ and $\GN$ makes the two $d$-choices
cancel; otherwise the choices contribute $1-1/p$.
After the factor $\mu(\ell)/\ell$, the tail $\ell>16U_0$ is
$O(\sum_{\ell>16U_0}\ell^{-2})=O(1/U_0)$.
The integer-counting errors total
$O(2^{\omega_{\rm all}(\GN)}\log U_0)=O(x^{o(1)})$.
The resulting $u$ main term is
\begin{equation}\label{eq:u-main}
 \frac{U_0|I_1|}{S\zeta(2)}
 \prod_{p\mid S}(1-p^{-2})^{-1}
 \prod_{p\mid\GN}(1+1/p)^{-1}.
\end{equation}
Multiplication by the $v$-sieve factor and
$|\operatorname{SL}_2(\Z/S\Z)|=S^3\prod_{p\mid S}(1-p^{-2})$
gives \eqref{eq:R-matrix} with \eqref{eq:cN}.
At $\GN=2$ this is $1/(3\zeta(2))$.
For completeness, split the factors $1-1/p$ at $p=L$.
Mertens' product estimate gives
$\prod_{p\mid\GN,p\le L}(1-1/p)\ge\prod_{p\le L}(1-1/p)
\gg1/\log L$.
For $p>L\ge2$, $\log(1-1/p)\ge-2/p$ and
$\sum_{p\mid\GN,p>L}1/p\le\omega_{\rm all}(\GN)/L\le1/\log2$.
Thus the product over these larger primes is bounded below by an
absolute positive constant. It follows from the squared-factor
formula in \eqref{eq:cN} that
\begin{equation}\label{eq:cN-lower}
 c_{\GN}\ge\zeta(2)^{-1}\prod_{p\mid\GN}(1-1/p)^2
      \gg(\log L)^{-2}=(\log\log\GN)^{-2}.
\end{equation}
The upper bound $c_{\GN}\le1$ is immediate.
Only that upper bound is needed for the root moment.

\emph{Error for $U_0\le V_0$.}
Summing the rectangle errors over $u\asymp U_0$ gives
\begin{equation}\label{eq:R-error}
 \begin{split}
 O\bigl(x^{o(1)}S^2(V_0U_0^{0.99}+U_0^{1.99})\bigr)
 &=O\bigl(U_0V_0x^{o(1)}S^2
                (U_0^{-0.01}+U_0^{0.99}/V_0)\bigr)\\
 &=O(U_0V_0x^{-c_0\delta}),\qquad c_0=1/200.
 \end{split}
\end{equation}
Both terms in parentheses are at most $x^{-\delta/100}$ and
$S^2\le\exp(2C_0L^{0.98})=x^{o(1)}$.
This fixes $c_0$ before the next subdivision.

\emph{The case $V_0<U_0$.}
Complete $(v,u)^T$ in the normalized way:
$c''u\equiv-\GN\pmod v$, $0\le c''<v$,
$d'''=(\GN+c''u)/v$. Put $\varepsilon=\GN/(uv)$.
Our hypothesis gives
\begin{equation}\label{eq:epsilon}
 0<\varepsilon\le\frac{\GN}{d_0YMN}<1,
 \qquad\frac{\GN}{d_0YMN}\ll\frac{L^B}{d_0Y}
 \quad\hbox{in the application}.
\end{equation}
Since $d''/v=\beta_P+\varepsilon$, $k=\lfloor d''/v\rfloor\in\{0,1\}$,
and
\begin{equation}\label{eq:swapped}
 \begin{split}
 c''&=(1+k)v-d'',\quad d'''=(1+k)u-c',\\
 \widetilde g&=\begin{pmatrix}v&(1+k)v-d''\\u&(1+k)u-c'\end{pmatrix},
 \quad\det\widetilde g=\GN,\\
 \beta''_P&=1+k-\beta_P-\varepsilon,\qquad
 k=1\ \Longleftrightarrow\ \beta_P+\varepsilon\ge1.
 \end{split}
\end{equation}
Either endpoint $\beta''_P\in\{0,1\}$ would imply $v\mid\GN$, impossible for
$v>1$ and $(v,\GN)=1$. Both $k$ branches must be included.
For $k=\kappa\in\{0,1\}$ the prescribed residue becomes
\begin{equation}\label{eq:swapped-residue}
 \widetilde g_0^{(\kappa)}=
 \begin{pmatrix}v_0&(1+\kappa)v_0-d''_0\\
                 u_0&(1+\kappa)u_0-c'_0\end{pmatrix},
 \quad
 J_\kappa(u,v)=\bigl((1+\kappa)-I_3-\GN/(uv)\bigr)\cap[0,1).
\end{equation}
Their union is $1-(I_3+\GN/(uv))$ modulo $1$, of length $|I_3|$.

Partition $I_1\times I_2$ into cells of relative side
$x^{-\kappa_0}$, $\kappa_0=c_0\delta/4$.
There are $O(x^{c_0\delta/2})$ cells. On each,
$\GN/(uv)$ varies by at most $4x^{-\kappa_0}$.
For each branch start with the affine interval
$(1+\kappa)-I_3-\varepsilon_0$, where $\varepsilon_0$ is its
corner value. Enlarge or contract this interval by $4x^{-\kappa_0}$
at each end and only then intersect with $[0,1)$, taking reversed
intervals as empty. This also covers a branch whose clipped
interval is empty at the corner.
Apply the shorter-coordinate case on that cell, with the two
coordinate roles interchanged and residue
$\widetilde g_0^{(\kappa)}$. Its error is independent of interval
lengths, so the summed error is
$O(U_0V_0x^{-c_0\delta/2})$.
The boundary mass is at most $8x^{-\kappa_0}$ per cell times its
main-term density, and totals $O(U_0V_0x^{-\kappa_0})$.
The two branches preserve $|I_3|$, giving final exponent
$c_0/4=1/800$ in \eqref{eq:R-matrix}.
This argument counts $\beta=c'/u$ directly. It does not differentiate
the map $r\mapsto\GN r$.

\emph{Prescribed lines.}
At $p\mid S$, matrices of determinant $\GN$ with a prescribed
first-row kernel number $p(p-1)$: there are $p-1$ nonzero first rows
with that kernel and $p$ choices of second row. The total number
is $p(p^2-1)$, so the proportion is $1/(p+1)$.
CRT multiplies these proportions. Summing the per-matrix errors
costs at most $S^3=x^{o(1)}$; decrease the fixed positive exponent
$c$ in \eqref{eq:R-lines}, for example to $1/1600$ for sufficiently
large $x$ with fixed $\delta$.
The factor at $2$ is already included in $c_{\GN}$ and is not
inserted again.
\end{proof}

\section{Transport of the Type II argument}\label{sec:transport}
\begin{lemma}[Reduced determinant and boundary geometry]\label{lem:transport-geometry}
For $Q=\gN(P)z$, $Q'=\gN(P)z'$ in one labelled block,
\begin{equation}\label{eq:reduced-det}
 \det(\gN(P)z,\gN(P)z')=\GN\det(z,z'),\qquad
 t'=\frac{\det(Q,Q')}{\GN D}=\frac{\det(z,z')}{D}.
\end{equation}
The kernel in relative coordinates is \eqref{eq:reduced-brace} with
cutoff $\psi_1(\det(z,z')/(DY))$. On allowed positions and edges,
\begin{equation}\label{eq:path-geometry}
 \frac1{16}\le\tauN(z)\le16,\qquad
 |\det(z,z')|\ll d_0Y,\qquad |z|\ll d_0YMN/\GN\le d_0Y.
\end{equation}
For $s=u/(d_0YM)$, $t_0=v/N$, the position boundaries are
\begin{equation}\label{eq:boundaries}
 s\tauN(z)\in\{1,16\},\qquad
 E_2=t_0\tauN(z)+\frac{z_2\GN}{d_0YMNs}\in\{1,2\}.
\end{equation}
The second term is bounded by $64$ in absolute value on allowed
positions; its $s$ derivative is $O(1)$, and
$\partial_{t_0}E_2=\tauN(z)\ge1/16$.
\end{lemma}
\begin{proof}
The determinant identities follow from $\det\gN=\GN$ and
\eqref{eq:positions}. The remaining bounds follow from
Lemma~\ref{lem:coordinates} and the cutoff in \eqref{eq:K}.
In the last boundary formula $s\in[1,16]$, so the extra term and its
$s$ derivative are $O(1)$; its $\beta$ derivative is
$t_0z_2=\exp(O(L^{0.21}))$. Transversality, rather than smallness of
this term, is what the cube count uses.
\end{proof}

The full dictionary is Table~\ref{tab:dictionary} in
Appendix~\ref{app:dictionary}. It includes the all-integer square,
the coefficient conjugation and the unit-coordinate support.
The position cutoffs in LWW Sections~4.2--4.7 are sharp;
$\eta$ is a cutoff in the list product. There is no smooth cutoff
in the second position coordinate.

\begin{proposition}[Root averaging]\label{prop:root-average}
Using the arguments in $\mathrm H2'$, the block version of
LWW Proposition~4.4 \LWW{pp.~22--24} is
\begin{equation}\label{eq:root-average}
 \begin{split}
 &\frac1{d_0YMN}\sum_{P\in\Omega\cap\CC_{\GN}}
       \langle u_P,(\mathcal A_{\GN}\mathcal A_{\GN}^*)^Ru_P\rangle\\
 &\quad=c_{\GN}\int_1^{16}\int_1^2\int_0^1
 M_R^{(\GN)}(d_0YMs,Nt_0,\beta)\,d\beta\,dt_0\,ds
       +O_A(L^{-2AR}).
 \end{split}
\end{equation}
Here $M_R^{(\GN)}$ is LWW's $M_R$ with $g'$ replaced by $\gN$
and LWW's coordinate-parity test replaced by the labelled state
space of Section~\ref{subsec:enlargement}.
\end{proposition}
\begin{proof}
A closed path from a root $P\in\CC_{\GN}$ has
$P_j=\gN(P)z_j$, with $z_j$ primitive,
$z_0=z_{2R}=(1,0)^T$ and $|z_j|\le2^7d_0Y$.
List divisibility becomes
$[z_j]_p=\ell_p(P)=\ker(u,c')\bmod p$.
Its good-state ratio is \eqref{eq:ratio-relative} at the current
state. Apply the path expansion and cube argument of LWW
Proposition~4.4 with Lemma~\ref{lem:R} in place of its root count.
The measure (4.29), p.~23, becomes
$c_{\GN}\,d(u/(d_0YM))\,d(v/N)\,d\beta$, and each prescribed line
has mass $1/(p+1)$.

The absolute coefficient mass remains $\exp(O_K(L^{0.73}))$,
since $|z_j|\ll Rd_0Y$.
The Bonferroni truncation has
$S\le\exp(O(L^{0.96}))\le\exp(C_0L^{0.98})$.
By \eqref{eq:boundaries} and
\begin{equation}\label{eq:ratio-derivative}
 \partial_\beta\frac{\tauN(w)}{\tauN(z)}
     =\frac{\det(z,w)}{\tauN(z)^2}
     =\frac1{\tauN(z)^2}\quad(\det(z,w)=1),
\end{equation}
the boundary surfaces have a transverse derivative at least $1/16$
and gradients $x^{o(1)}$. The neighborhood volume is
$\ll\varepsilon\max|\nabla E|/|\partial_{t_0}E|
=x^{-c\delta/8+o(1)}$ at the chosen cube size.
Sum the per-cube power errors before the moment limit; their
product with the coefficient mass gives $O_A(L^{-2AR})$.
These are the changes to the stated LWW argument.
\end{proof}

\begin{proposition}[Block Type II estimate]\label{prop:II}
Assume $\mathrm H2'$. Let $M,N\ge x^\delta$, $MN\asymp x$ and
let the coefficients satisfy the sifting and bound hypotheses above.
After choosing $A_0$ and then $K$ sufficiently large, independently
of the fixed band exponents $a_i$,
\begin{equation}\label{eq:II}
 |\QQ-\Qmaj|\ll_A\GN YL^{-A}.
\end{equation}
\end{proposition}
\begin{proof}
We apply the arguments of LWW Sections~4.4--4.7 as written, with
the replacements in Table~\ref{tab:dictionary}.
Section~4.4 uses relative coordinates, the linear form $\tau$ and
the projection $\chi$; use $\tauN$ and
$\chi_{\GN}=\ind_{\gN z\in\Omega}$.
Its auxiliary operators have the same relative determinant,
shared-prime lines and intrinsic tests. In (4.40)--(4.41),
\LWW{pp.~27--28},
\begin{equation}\label{eq:bezout-edge}
 z'=j_1w_z+hz,\qquad
 \frac{\tauN(z')}{\tauN(z)}
     =h+j_1\frac{\tauN(w_z)}{\tauN(z)}.
\end{equation}
The ratio strip leaves $O(1)$ integers $h$, uniformly in $j_1$.
The clean and dirty bounds (4.42)--(4.50) use this geometry and the
same bounded endpoint multipliers.
The combinatorial signed-memory expansion of Section~4.5,
pp.~32--35, uses these operator bounds, root averaging and bounded
measure mass; $c_{\GN}\le1$ suffices.

For the endpoint pairing in Section~4.6, use the inner product
conjugate-linear in its initial entry and set
\begin{equation}\label{eq:endpoint-conjugation}
 f(Db m,s)=\overline{\alpha_m}\beta_s.
\end{equation}
At $P=(Dbm,s)$, $Q=(Dar,n)$ this gives exactly
\begin{equation}\label{eq:pairing-product}
 \overline{f(P)}f(Q)
 =\alpha_m\overline{\alpha_r}\beta_n\overline{\beta_s}.
\end{equation}
All support, list-invariance and norm bounds are unchanged by this
conjugation. Every prime divisor of $Db$ is below $Z$, whereas $s$
is $Z$-rough; a common prime of the endpoint coordinates must therefore
divide $m,s$. If $p\mid(m,s)$ on endpoint support, then $p>Z$ and
$p\nmid\GN$; the determinant equation gives $p\mid b-a$.
The retained cross-edge exclusion gives $b\ne a$, while
$0<|b-a|<4Y<Z$, a contradiction. The argument is independent of
the sign of the common cofactor.

Use Lemma~\ref{lem:block-size} for the partition in Section~4.6,
and the positive matrix and ambient spectral H\"older argument in
Lemma~\ref{lem:compression} for (4.64)--(4.68), pp.~36--37.
The trace loss is $(d_0Y)^{1/(2R)}=O(1)$.
Section~4.7 has $O(1)$ exact-edge and $O(Y)$ comparison-edge counts
in the reduced determinant, as in (4.70), p.~37.
For (4.75), p.~38, fix $k_0=mn$, $\ell_0=rs$.
Then
\begin{equation}\label{eq:comparison-count}
 \ell_0\equiv ba^{-1}k_0\pmod{\GN},\qquad
 |bk_0-a\ell_0|\ll\GN Y.
\end{equation}
The interval for $\ell_0$ has length $O(\GN)$ and one residue
class, so there are $O(1)$ choices. Raw absolute estimates use the
all-integer cofactor range of Proposition~\ref{prop:cofactor-square};
the extra $\GN/X$ costs a fixed logarithmic power.

The moment/root normalization itself remains $d_0YMN$, giving the
pairing scale $XY$ of LWW. The right side $\GN Y$ in
\eqref{eq:II} is a permitted relaxation by
$\GN/X\ll L^B$, absorbed by increasing the fixed precision.
It is independent of the Cauchy cofactor factor $\GN/Y$.
\end{proof}
\begin{remark}[Determinant scale]\label{rem:det-scale}
Using arcs for the unreduced determinant without the block
congruence would connect a position to $\asymp\GN Y$ positions.
The trace loss would be
$(d_0Y\GN)^{1/(2R)}=\exp(L^{1/2}+o(L^{1/2}))$ rather than $O(1)$.
The block congruence absorbs this determinant scale and reappears
in the major arcs.
\end{remark}

\section{Major arcs}\label{sec:major}
\subsection{Reduction to the averaged energy}\label{subsec:arc-reduction}
\begin{proposition}[Major-arc reduction]\label{prop:arc-reduction}
For every fixed $D_*>0$, \eqref{eq:alpha-star} at a sufficiently
large fixed precision $A=D_*+O(A_0+B_1)$ implies
\begin{equation}\label{eq:major-bound}
 |\Qmaj|\ll_{D_*}\GN YL^{-D_*}.
\end{equation}
\end{proposition}
\begin{proof}
By Lemma~\ref{lem:cover}, on the six-variable unit support,
\begin{equation}\label{eq:character-congruence}
 \begin{split}
 H_{\GN}(bmn-ars;a,b)
 &=\psi_1(t'/Y)\ind_{bmn\equiv ars\ (\GN)}
      \int_{\Mone}e(\phi(t'-b+a))\,d\phi,\\
 t'&=(bmn-ars)/\GN,\qquad
 \ind_{bmn\equiv ars\ (\GN)}
 =\frac1{\ph(\GN)}\sum_{\chi\bmod\GN}\chi(bmn)\overline{\chi(ars)}.
 \end{split}
\end{equation}
Write $\phi=h_0/k+\lambda_0/Y$,
$(h_0,k)=1$, $k\le L^{A_0}$, $|\lambda_0|\le2L^{A_0}$,
and $d\phi=d\lambda_0/Y$. Its factor is
\begin{equation}\label{eq:arc-phases}
 e(h_0t'/k)e(-h_0(b-a)/k)
 e(\lambda_0t'/Y)e(-\lambda_0(b-a)/Y).
\end{equation}
The factors in $a,b$ modulo $k$ expand in characters, with mass
at most $k$ each, since $a,b$ are units modulo $k$.

For $e(h_0t'/k)$, put $u_1=bmn$, $u'_1=ars$.
Both are units modulo $k\GN$, and
\[
 \ind_{\GN\mid u_1-u'_1}e(h_0t'/k)
 =\ind_{\GN\mid u_1-u'_1}
    e(h_0u_1/(k\GN))e(-h_0u'_1/(k\GN)).
\]
Let
\begin{equation}\label{eq:fibre-group}
 H=\ker((\Z/k\GN\Z)^\times\longrightarrow(\Z/\GN\Z)^\times),
 \qquad |H|=k\prod_{\substack{p\mid k\\p\nmid\GN}}(1-1/p)\le k.
\end{equation}
The reduction is surjective. For each unit class $a'\bmod\GN$,
choose a unit representative $\widetilde a'\bmod k\GN$; its fibre
is $\widetilde a'H$. Expand
$w\mapsto e(h_0\widetilde a'w/(k\GN))$ in characters
$\eta\in\widehat H$, with coefficients $|c_{a',\eta}|\le1$.
Every such character extends to a character $\psi_\eta$ of the
finite abelian group $(\Z/k\GN\Z)^\times$, independently of $a'$.
Hence on that fibre
\begin{equation}\label{eq:fibre-expansion}
 e(h_0u_1/(k\GN))
 =\sum_{\eta\in\widehat H}
 c_{a',\eta}\psi_\eta(\widetilde a')^{-1}\psi_\eta(u_1).
\end{equation}
The same representative works on the $u'_1$ side. Cauchy gives
\begin{equation}\label{eq:fibre-cauchy}
 \sum_{a'}\left|\sum_\eta
 c_{a',\eta}\psi_\eta(\widetilde a')^{-1}X_{a',\eta}\right|^2
 \le k\sum_\eta\sum_{a'}|X_{a',\eta}|^2.
\end{equation}
Thus class-dependent coefficients disappear at fixed powers of
$k$, including when $(k,\GN)>1$. Together with the $a,b$ factors,
the character mass is at most $k^4$ (LWW use $k^6$).

For the smooth factor put
$w=u_1/(\GN Y)$, $w'=u'_1/(\GN Y)$ and
$\kappa_X=MN/\GN$. The dyadic variables satisfy
\begin{equation}\label{eq:natural-scale}
 w,w'\in[\kappa_X,16\kappa_X],\qquad
 \kappa_X\gg L^{-B},\quad\kappa_X\le1,
 \qquad
 \psi_1(t'/Y)e(\lambda_0t'/Y)
       =\psi_1(w-w')e(\lambda_0(w-w')).
\end{equation}
If $\kappa_X\le1/4$, the difference has absolute value at most
$4$ and $\psi_1=1$. Otherwise it varies on a bounded interval.
Set $z=\log(w'/\kappa_X)$ and $\zeta=\log(w/w')$.
With a fixed smooth cutoff $\varpi$ equal to $1$ on $[1,16]$,
the kernel is
\begin{equation}\label{eq:mellin-kernel}
 \varpi(e^z)\varpi(e^{z+\zeta})
 \psi_1(\kappa_Xe^z(e^\zeta-1))
 e(\lambda_0\kappa_Xe^z(e^\zeta-1)).
\end{equation}
Its support lies in a fixed box of side $O(1)$.
In the variables $w/\kappa_X,w'/\kappa_X$, derivatives are bounded
by fixed powers of $L^{A_0}$ because $\kappa_X\le1$.
Consequently for a fixed sufficiently large $B_0$ the transform obeys
\begin{equation}\label{eq:mellin-decay}
 F_{\lambda_0}(v,s)\ll_j
 (1+|v|/L^{B_0})^{-j}(1+|s|/L^{B_0})^{-j}.
\end{equation}
Taking $B_1$ sufficiently large leaves frequencies
$|t|,|s|\le L^{B_1}$. This is the Mellin separation of LWW
Section~5, pp.~42--43, at scale $1$ rather than $X$.
The cutoff here does not force a narrow product difference;
the arithmetic is in the character congruence.

With the normalized polynomials \eqref{eq:polynomials}, each product
sum has size $YMN\cdot BMN$ with its displayed normalization.
The $1/Y$ arc measure therefore gives the scale
\begin{equation}\label{eq:major-normalization}
 \frac{(YMN)^2}{Y}\,\frac1{\ph(\GN)}
       \sum_\chi\int|BMN|^2
 =Y(MN)^2\,\frac1{\ph(\GN)}\sum_\chi\int|BMN|^2.
\end{equation}
To obtain \eqref{eq:major-bound} it suffices that the energy be
$\ll L^{-D_*}\GN/(MN)^2$, which follows from
$L^{-A}/\GN$ because $MN\le\GN$.
Cauchy in $t,\chi$, and the dyadic tail estimates of LWW
Section~5, p.~43, leave \eqref{eq:alpha-star} with
$\xi=\chi\psi$ on $m,n$ and $\xi_b=\chi\psi_2$ on $b$,
where $\psi=\psi_\eta$, $\psi_2=\psi\psi'_b$ and
$\psi'_b$ is modulo $k$. The primed product uses
$-\lambda_0,\overline\alpha,\overline\beta$; the estimates below
also hold after these conjugations and the corresponding sign
change of the frequency twist.
The arc count $L^{2A_0}$, character mass $k^4$, $\lambda_0$ integral
$L^{A_0}$ and $s$ integral $L^{B_1}$ are fixed logarithmic costs,
absorbed by $A=D_*+O(A_0+B_1)$.

For completeness, Parseval for a product polynomial gives
\begin{equation}\label{eq:trivial-energy}
 \frac1{\ph(\GN)}\sum_\chi
 \left|\sum_u c_u\chi(u)\psi(u)u^{it}\right|^2
 =\sum_{a'\in(\Z/\GN\Z)^\times}
      \left|\sum_{u\equiv a'\ (\GN)}c_u\psi(u)u^{it}\right|^2
 \ll(YX/\GN+1)\sum_u|c_u|^2
 \ll\frac{L^{O(1)}}{\GN}.
\end{equation}
Here the product length is $\asymp XY$ and
$|c_u|\ll L^{C}\tau_3(u)/(XY)$.
For $MN$ alone the squared coefficient norm is
$L^{O(1)}/(MN)\le L^{O(1)+B}/\GN$; this extra power is absorbed.
\end{proof}

\subsection{Prime sums from a fixed zero-free region}\label{subsec:P}
\begin{lemma}[Lemma P]\label{lem:P}
Let $2\le q\le\GN^3$, and let $\chi\bmod q$ be induced by a
primitive $\chi^*$ of conductor $q^*>1$.
Fix positive $C_{\mathrm{zf},1},C_{\mathrm{zf},2}$ and
$C_3\ge1$ with $C_{\mathrm{zf},2}\ge C_3+1$.
Suppose $L(s,\chi^*)\ne0$ in
\eqref{eq:Hstar}'s rectangle for these constants.
For $|u|\le L^{C_3}$ and $\GN^{1/6}\le y\le\GN$,
\begin{equation}\label{eq:P}
 S(y):=\sum_{\ell\le y}\Lambda(\ell)\chi(\ell)\ell^{iu}
 \ll yL^3(L^{-C_{\mathrm{zf},1}/12}+L^{-C_{\mathrm{zf},2}}).
\end{equation}
The implied constant can be absolute once
$L\ge L_0(C_{\mathrm{zf},1},C_{\mathrm{zf},2},C_3)$.
\end{lemma}
\begin{proof}
The imprimitive difference is supported on prime powers at primes
dividing $q$, and contributes at most
$\omega_{\rm all}(q)\log y\le5L^2$, absorbed since
$y\ge\GN^{1/6}$.
For $\chi^*$ put
\begin{equation}\label{eq:Perron-data}
 F(s)=-\frac{L'}{L}(s-iu,\chi^*),\quad
 c=1+1/L,\quad
 T=L^{C_{\mathrm{zf},2}}-L^{C_3}-1\ge\tfrac12L^{C_{\mathrm{zf},2}}.
\end{equation}
Truncated Perron \cite[Chapter~17]{Davenport} gives
\begin{equation}\label{eq:Perron}
 \sum_{\ell\le y}\Lambda(\ell)\chi^*(\ell)\ell^{iu}
 =\frac1{2\pi i}\int_{c-iT}^{c+iT}F(s)\frac{y^s}{s}\,ds
       +O(y(\log y)^2/T+\log y).
\end{equation}
The usual error is the sum of
$\Lambda(\ell)(y/\ell)^c
\min(1,(T|\log(y/\ell)|)^{-1})$.
The choice $y\le\GN$ ensures $y^c\le ey$ and
$\sum\Lambda(\ell)\ell^{-c}\ll L$.

For $1/2\le\sigma\le2$, the partial-fraction formula and local
zero count \cite[Chapter~16]{Davenport} read
\begin{equation}\label{eq:partial-fraction}
 \frac{L'}L(s,\chi^*)
 =\sum_{|\gamma-t|\le1}\frac1{s-\rho_0}
       +O(\log(q^*(|t|+2))),\quad
 N(t+1,\chi^*)-N(t,\chi^*)\ll\log(q^*(|t|+2)).
\end{equation}
Put $\sigma_0=1-C_{\mathrm{zf},1}\log L/(2L)$.
For $\sigma\ge\sigma_0$ and $|t|+1\le L^{C_{\mathrm{zf},2}}$,
all nearby zeros have real part at most
$1-C_{\mathrm{zf},1}\log L/L$ and horizontal distance at least
$C_{\mathrm{zf},1}\log L/(2L)$.
Their number is $O(\log(q^*L^{C_{\mathrm{zf},2}}))=O(L)$ after a
threshold depending on the fixed constants. Hence
\begin{equation}\label{eq:log-derivative}
 |L'/L(s,\chi^*)|\ll
 L\frac{2L}{C_{\mathrm{zf},1}\log L}+L\ll L^2.
\end{equation}
The shifted $F(s)$ has this bound on $\sigma\ge\sigma_0$, $|t|\le T$.

Move the contour to $\sigma_0$.
There is no zero pole in the rectangle: a pole $s=\rho_0+iu$ would
have $\Re\rho_0\ge\sigma_0$ and
$|\Im\rho_0|\le T+|u|\le L^{C_{\mathrm{zf},2}}-1$, contradicting
the region. Also $\sigma_0>1/2$, so $s=0$ is outside, and
$\chi^*$ is nonprincipal. The vertical and horizontal integrals
are respectively
\[
 \ll y^{\sigma_0}L^2\int_{-T}^T\frac{dt}{|\sigma_0+it|}
 \ll y^{\sigma_0}L^3,
 \qquad
 \ll(c-\sigma_0)y^cL^2/T\ll yL^2/T.
\]
Since $L/6\le\log y\le L$,
$y^{\sigma_0}\le yL^{-C_{\mathrm{zf},1}/12}$.
Combining these estimates yields
\begin{equation}\label{eq:P-end}
 S(y)\ll yL^{3-C_{\mathrm{zf},1}/12}+yL^{2-C_{\mathrm{zf},2}}
       +y(\log y)^2/T+\log y+L^2
 \ll yL^3(L^{-C_{\mathrm{zf},1}/12}+L^{-C_{\mathrm{zf},2}}).
\end{equation}
\end{proof}
\begin{remark}\label{rem:width}
The width $C\log L/L$ is what gives the arbitrary fixed logarithmic
saving on $y\ge\GN^{1/6}$. A width $c/\log q^*$ with
$q^*\asymp\GN$ gives only a factor of the size $e^{-c/6}$.
Only the logarithmic derivative on the contour is used.
\end{remark}

\begin{corollary}[Pointwise prime-factor polynomial]\label{cor:P}
Suppose $\xi=\chi\psi\bmod k\GN$ is induced by $\chi^*$ of
conductor $L^{A_1}<q^*\le k\GN$ and Lemma~\ref{lem:P} applies.
If
\begin{equation}\label{eq:P-constants}
 C_{\mathrm{zf},1}\ge12(D+4),\quad
 C_{\mathrm{zf},2}\ge D+4,\quad
 C_{\mathrm{zf},2}\ge C_3+1,\quad C_3=\max(B_1,2B)+1,
\end{equation}
then, uniformly for $|t|\le T_0+1$ and sufficiently large $L$
depending on the fixed parameters,
\begin{equation}\label{eq:N-pointwise}
 |N_\xi(t)|\le8L^{1-D}+\frac{4Z+2L}{V}\le9L^{1-D}.
\end{equation}
\end{corollary}
\begin{proof}
Put $u=t+v$, so $|u|\le L^{C_3}$.
Both $R_Z$ and the unit indicator are completely multiplicative.
Thus
\begin{equation}\label{eq:N-convolution}
 \begin{split}
 c_V(n)R_Z(n)\ind_{(n,\GN)=1}
 &=\sum_{\substack{\ell k'=n\\\ell>V}}
 \Lambda(\ell)R_Z(\ell)\ind_{(\ell,\GN)=1}
 R_Z(k')\ind_{(k',\GN)=1},\\
 N_\xi(t)&=\frac1N\sum_{k'<2N/V}
 R_Z(k')\ind_{(k',\GN)=1}\xi(k')(k')^{iu}
 \sum_{\ell\in I_{k'}}\Lambda(\ell)R_Z(\ell)
                    \ind_{(\ell,\GN)=1}\xi(\ell)\ell^{iu},\\
 I_{k'}&=(\max(V,N/k'),2N/k'].
 \end{split}
\end{equation}
In a ratio-two interval at most one power of each prime is present.
Dropping the indicators in the inner sum therefore costs at most
$\sum_{p\le Z}\log p+\sum_{p\mid\GN}\log p\le2Z+L$,
and after the $k'$ sum at most $(4Z+2L)/V$.
For every nonempty $I_{k'}$, the two endpoints are between $V$
and $2N\le\GN$, and
\[
 \log V=(\log x)/5\ge(L-B\log L-\log2)/5\ge L/6
\]
for sufficiently large $L$. The modulus is at most $\GN^2$.
Lemma~\ref{lem:P} and \eqref{eq:P-constants} give
$|S(y)|\le2C_PyL^{-D-1}\le yL^{-D}$ for $L\ge2C_P$.
Consequently
\[
 |N_\xi(t)|\le4L^{-D}(\log(2N/V)+1)+(4Z+2L)/V
 \le8L^{1-D}+(4Z+2L)/V.
\]
Since $V\ge\GN^{1/6}$ and $Z=\exp(L^{0.24})$, the last ratio is
at most $L^{1-D}$ for large $L$. Only endpoint sums with $y\ge V$
are used, even when $I_{k'}$ is short.
\end{proof}

The mean value theorem used below has the form
\begin{equation}\label{eq:MV-general}
 \int_{-T}^T\left|\sum_u b_u u^{it}\right|^2dt
 =\sum_u|b_u|^2\bigl(2T+O(\delta_u^{-1})\bigr),
\end{equation}
where $\delta_u$ is the separation of $\log u$ from the other
logarithmic frequencies \cite{MV}.

\subsection{Exceptional values of the short polynomial}\label{subsec:LS}
For $0\le i<2T_0$ let $J_i=[-T_0+i,-T_0+i+1]$.
For each $\chi\bmod\GN$ choose a point $t_{\chi,i}$ maximizing
$|B_{\xi_b}M_\xi N_\xi|$ on $J_i$.
Let $\mathcal P_0$ use the even $i$ and $\mathcal P_1$ the odd $i$.
Within each family, points for a fixed character are $1$-spaced, and
\begin{equation}\label{eq:discretize}
 \sum_\chi\int_{|t|\le T_0}|BMN|^2\,dt
 \le\sum_{\varepsilon=0,1}
       \sum_{(\chi,t)\in\mathcal P_\varepsilon}
        |B_{\xi_b}(t)M_\xi(t)N_\xi(t)|^2.
\end{equation}
Below $\mathcal P$ is either family, and each argument applies to both.

\begin{lemma}[Sobolev estimate]\label{lem:Sobolev}
For $F\in C^1([-T_0-1,T_0+1])$ and $1$-spaced points
$t_1<\cdots<t_r$ in $[-T_0,T_0]$,
\begin{equation}\label{eq:Sobolev}
 \sum_j|F(t_j)|^2\le\int_{|t|\le T_0+1}|F|^2\,dt
 +2\left(\int_{|t|\le T_0+1}|F|^2\,dt\right)^{1/2}
   \left(\int_{|t|\le T_0+1}|F'|^2\,dt\right)^{1/2}.
\end{equation}
\end{lemma}
\begin{proof}
Average
$|F(t_j)|^2=|F(s)|^2+\int_s^{t_j}(|F|^2)'$
over $s\in[t_j-1/2,t_j+1/2]$.
These intervals have disjoint interiors; use
$|(|F|^2)'|\le2|F||F'|$ and Cauchy--Schwarz.
\end{proof}

\begin{lemma}[Mean value over residue classes]\label{lem:MV}
If $(c_u)$ is supported on positive $u\le4\GN$ and $|\psi_3|\le1$,
then for $T\ge1$,
\begin{equation}\label{eq:MV}
 \sum_{\chi\bmod\GN}\int_{-T}^T
 \left|\sum_uc_u\chi(u)\psi_3(u)u^{it}\right|^2dt
 \le\ph(\GN)(2T+C_0')\sum_u|c_u|^2
\end{equation}
with absolute $C_0'$.
\end{lemma}
\begin{proof}
Character orthogonality gives
\[
 \sum_\chi\left|\sum_uc_u\chi(u)d_u\right|^2
 =\ph(\GN)\sum_{a'\in(\Z/\GN\Z)^\times}
           \left|\sum_{u\equiv a'\ (\GN)}c_ud_u\right|^2.
\]
Each class has at most four positive integers at most $4\GN$.
If $u<u'$ are in it, $u'\ge u+\GN\ge(4/3)u$ since
$u\le3\GN$. The Montgomery--Vaughan mean value theorem
\cite{MV} therefore gives $C_0'=O(1/\log(4/3))$.
For length at most $\GN$, a unit class has at most one term and the
integral in that class is exactly $2T$ times its squared coefficient.
\end{proof}

\begin{lemma}[Lemma LS]\label{lem:LS}
Fix $A_s\ge1$ and put
$E=\{(\chi,t)\in\mathcal P:|B_{\xi_b}(t)|>L^{-A_s}\}$.
For $L\ge L_0(K,a_K,A_s,B_1,C_B)$,
\begin{equation}\label{eq:LS}
 |E|\le\exp\bigl(L^{1-a_K}(2A_s+K+10)\log L\bigr)
       \le\exp(L^{0.9}).
\end{equation}
\end{lemma}
\begin{proof}
Let $j=\lfloor L/\log(16Y)\rfloor\ge1$; then
$(16Y)^j\le\GN<(16Y)^{j+1}$.
Write $B_{\xi_b}(t)^j=\sum_uc_u^{(j)}\xi_b(u)u^{it}$, where
\begin{equation}\label{eq:B-power}
 c_u^{(j)}=\sum_{\substack{b_1\cdots b_j=u\\b_i\in\Aset}}
      \prod_i c_{b_i},\qquad
 c_b=\frac{\eta(b/Y)e(\lambda_0b/Y)}{Y\prod_iV_i},\qquad
 |c_b|\le C_B/(16Y).
\end{equation}
This polynomial is supported on $[Y^j,(16Y)^j]\subset[1,\GN]$.
For each disjoint band the $j$ prime factors of $u$, with
multiplicity, can be assigned to the $j$ slots in at most $j!$ ways.
Hence there are at most $(j!)^K$ representations and
\begin{equation}\label{eq:B-power-norm}
 |c_u^{(j)}|^2\le(j!)^K\sum_{\rm reps}\prod_i|c_{b_i}|^2,
 \qquad
 \sum_u|c_u^{(j)}|^2
 \le(j!)^K\left(\sum_b|c_b|^2\right)^j
 \le(j!)^K\left(\frac{C_B^2}{16Y}\right)^j.
\end{equation}
Use Lemmas~\ref{lem:Sobolev} and~\ref{lem:MV} on $B^j$;
its derivative inserts a coefficient $i\log u$, with $|\log u|\le L$.
Since $\xi_b=\chi\psi_2$ and $|\psi_2|\le1$,
\begin{equation}\label{eq:B-moment}
 \sum_{(\chi,t)\in\mathcal P}|B_{\xi_b}(t)|^{2j}
 \le3L\ph(\GN)(2T_0+2+C_0')(j!)^K(C_B^2/(16Y))^j.
\end{equation}
The left side is at least $|E|L^{-2jA_s}$.
Since $Y^{-j}<16^{j+1}Y/\GN$ and $\ph(\GN)\le\GN$,
\begin{equation}\label{eq:E-count}
 |E|\le L^{2jA_s}3L(2T_0+2+C_0')(j!)^KC_B^{2j}16Y.
\end{equation}
Every $b\in\Aset$ exceeds $\exp(L^{a_K})$, so
$j\le L^{1-a_K}$, $\log j\le\log L$, and
$\log Y\le2\sum_iL^{a_i}+\log16\le3KL^{a_K}$.
It follows that
\[
 \log|E|\le2A_sj\log L+Kj\log j+2j\log C_B+\log Y+O(\log L)
 \le L^{1-a_K}(2A_s+K+10)\log L.
\]
For fixed $a_K>1/10$, this is eventually at most $L^{0.9}/2$.
\end{proof}

\begin{corollary}[Bulk estimate]\label{cor:bulk}
For each family $\mathcal P$,
\begin{equation}\label{eq:bulk}
 \begin{split}
 \sum_{\mathcal P\setminus E}|BMN|^2
 &\le L^{-2A_s}\sum_\mathcal P|M_\xi N_\xi|^2\\
 &\le L^{-2A_s}3L\ph(\GN)(2T_0+O(1))
           \frac{L^2\sum_{u\le4MN}\tau(u)^2}{(MN)^2}\\
 &\ll L^{-2A_s+B_1+B+8}\frac{\ph(\GN)}{\GN}.
 \end{split}
\end{equation}
\end{corollary}
\begin{proof}
The product $MN$ polynomial is supported in $(MN,4MN]$, with
$4MN\le8x\le4\GN$ and coefficients at most $L\tau(u)/(MN)$.
Apply Lemmas~\ref{lem:Sobolev} and~\ref{lem:MV}, and
$\sum_{u\le y}\tau(u)^2\ll y(\log y)^3$.
\end{proof}

\subsection{Energy on the exceptional set}\label{subsec:HM}
\begin{lemma}[Lemma HM]\label{lem:HM}
Let $(\xi_r,t_r)_{r\le R'}$ be character-frequency pairs modulo $q$,
$1$-spaced in $t$ whenever their characters agree.
Let $a_m$ be supported on $(X',X'']$, $X'<X''\le8X'$, and let
$w\ge0$ be at least $1$ on that interval. Then
\begin{equation}\label{eq:HM}
 \sum_r\left|\sum_ma_m\xi_r(m)m^{it_r}\right|^2
 \le\|a\|_2^2\max_r\sum_s|G_{rs}|,\qquad
 G_{rs}=\sum_mw(m)\xi_r(m)\overline{\xi_s(m)}m^{i(t_r-t_s)}.
\end{equation}
\end{lemma}
\begin{proof}
Put $S_r=\sum_ma_m\xi_r(m)m^{it_r}$ and
$b_r=\overline{S_r}/(\sum_s|S_s|^2)^{1/2}$; the zero case is
immediate. Cauchy gives
\[
 (\sum_r|S_r|^2)^{1/2}=\sum_rb_rS_r
 \le\|a\|_2\left(\sum_mw(m)
           \left|\sum_rb_r\xi_r(m)m^{it_r}\right|^2\right)^{1/2}.
\]
The squared last factor is $\sum_{r,s}b_r\bar b_sG_{rs}$.
The Hermitian Schur bound and $\sum|b_r|^2=1$ prove the claim.
\end{proof}

Choose $q=k\GN$ and a smooth $w_0$ with
$\ind_{(X',X'']}\le w_0\le\ind_{[X'/2,16X']}$ and
$|w_0^{(i)}|\le C_i(X')^{-i}$. There are three kinds of Gram entry.
\begin{enumerate}
\item On the diagonal, $G_{rr}\le16X'$.
\item If $\xi_r=\xi_s$, $r\ne s$, and $\tau=t_r-t_s$, then
$|\tau|\ge1$ and
\begin{equation}\label{eq:Gram-same}
 G_{rs}=\sum_{d\mid q}\mu(d)d^{i\tau}
              \sum_{m'}w_0(dm')(m')^{i\tau}.
\end{equation}
For $d\le16X'$, Euler summation replaces the inner sum by its
integral with error $O(1+|\tau|)$: for
$f(y)=w_0(dy)y^{i\tau}$,
\[
 \int|f'|\le\int|w_0'|
       +|\tau|\int_{X'/(2d)}^{16X'/d}dy/y\le C(1+|\tau|).
\]
Two integrations by parts give integral size
$O((X'/d)(1+|\tau|)^{-2})$.
For $d>16X'$ the sum is empty. Thus
\begin{equation}\label{eq:Gram-same-bound}
 |G_{rs}|\le C'X'(1+|\tau|)^{-2}\log L
                +\tau(k\GN)(1+|\tau|),
\end{equation}
using $\sum_{d\mid q}1/d\ll\log\log q$.
The $1$-spacing gives
\begin{equation}\label{eq:Gram-same-row}
 \sum_{\substack{s\ne r\\\xi_s=\xi_r}}|G_{rs}|
 \le C''X'\log L+(2T_0+1)^2\tau(k\GN).
\end{equation}
\item If $\xi_r\ne\xi_s$, their quotient is nonprincipal.
On units it is $\chi_r\overline{\chi_s}$, and these are distinct
characters modulo $\GN$.
Partial summation gives
\begin{equation}\label{eq:Gram-distinct}
 |G_{rs}|\le C(1+|\tau|)\Delta_\xi(16X'),\qquad
 \Delta_\xi(Y')=\max_{y\le Y'}\left|\sum_{m\le y}\xi(m)\right|,
 \quad |\tau|\le2T_0+1.
\end{equation}
If $\xi$ is induced by primitive $\xi^*\bmod q^*>1$, put
$q_1=\prod_{p\mid k\GN,p\nmid q^*}p$. Then
\begin{equation}\label{eq:imprimitive}
 \sum_{m\le y}\xi(m)
 =\sum_{d\mid q_1}\mu(d)\xi^*(d)\sum_{m'\le y/d}\xi^*(m'),
 \qquad \Delta_\xi(Y')\le2^{\omega_{\rm all}(k\GN)}\Delta_{\xi^*}(Y').
\end{equation}
The two inputs are P\'olya--Vinogradov and Burgess at $r=2$:
\begin{align}
 \Delta_{\xi^*}(Y')&\le C(q^*)^{1/2}\log q^*
       \le C(L^{A_0}\GN)^{1/2}\,2L,\label{eq:PV}\\
 \left|\sum_{H_0<n\le H_0+H}\xi^*(n)\right|
   &\ll_\varepsilon H^{1/2}(q^*)^{3/16+\varepsilon}
       \quad(H\ge1),\label{eq:Burgess}
 \end{align}
where \eqref{eq:Burgess} holds for every modulus, as in
\cite[introduction, (1)]{HeathBrown2012} (arbitrary moduli are
allowed for $r\le3$).
Also
\[
 2^{\omega_{\rm all}(k\GN)}
 \le2^{A_0\log L}2^{\omega_{\rm all}(\GN)}
 \le L^{A_0}\exp(O(L/\log L))=\GN^{o(1)},\qquad
 \tau(k\GN)=\GN^{o(1)}.
\]
\end{enumerate}

\begin{corollary}[Exceptional energy]\label{cor:exceptional-energy}
Let $E$ be as in Lemma~\ref{lem:LS} and fix $\eta_0=1/40$.
Then
\begin{align}
 M\ge\GN^{3/8+\eta_0}&\quad\Longrightarrow\quad
        \sum_E|M_\xi(t)|^2\ll\log L,\label{eq:EM}\\
 \sum_E|M_\xi(t)|^6&\ll L^9,\label{eq:EM3}\\
 N\ge\GN^{1/2+2\eta_0}&\quad\Longrightarrow\quad
        \sum_E|N_\xi(t)|^2\ll L^2\log L.\label{eq:EN}
\end{align}
\end{corollary}
\begin{proof}
For \eqref{eq:EM} use $a_m=\alpha_m/M$,
$\|a\|_2^2\le2/M$, $X'=M$ and Burgess with
$\varepsilon=\eta_0/8$.
The distinct-character row contribution is at most
\begin{equation}\label{eq:EM-margin}
 \begin{split}
 &|E|C(2T_0+2)2^{\omega_{\rm all}(k\GN)}
       C_\varepsilon(16M)^{1/2}(L^{A_0}\GN)^{3/16+\varepsilon}\\
 &\hspace{4mm}\le M\,\GN^{-\eta_0/2+\varepsilon+o(1)}
                      \exp(L^{0.9})
   =M\,\GN^{-3/320+o(1)}\exp(L^{0.9})=o(M).
 \end{split}
\end{equation}
Together with $16M+C''M\log L+(2T_0+1)^2\tau(k\GN)$ this proves
\eqref{eq:EM}.
For \eqref{eq:EM3},
$M^3\ge(x^{1/5}/2)^3\ge\GN^{3/5-o(1)}$.
The coefficients of $M_\xi^3$ are at most $M^{-3}\tau_3(u)$ on
$(M^3,8M^3]$, with squared norm $\ll L^8/M^3$.
P\'olya--Vinogradov makes the distinct-character row at most
\[
 |E|L^{O(1)}\GN^{1/2+o(1)}
 \le M^3\GN^{1/2-3/5+o(1)}\exp(L^{0.9})=o(M^3),
\]
yielding $L^9$.
For \eqref{eq:EN}, take $a_n=\beta_n/N$,
$\|a\|_2^2\le2L^2/N$; P\'olya--Vinogradov gives row
$\le N\GN^{-2\eta_0+o(1)}\exp(L^{0.9})=o(N)$.
In all cases the same-character row is $O(X'\log L)$, because
$X'\ge\GN^{1/5-o(1)}$ absorbs $(2T_0+1)^2\tau(k\GN)$.
\end{proof}

\begin{inputthm}[Twisted sifted M\"obius sums]\label{input:sifted}
LWW Proposition~2.2, (2.5), p.~5, states that for fixed
$\delta,B',C'>0$, $x^\delta\le D_1\le2x$,
$I\subset[D_1,2D_1]$ and a character of modulus
$q\le(\log x)^{B'}$, for every fixed $A'>0$,
\begin{equation}\label{eq:LWW-sifted}
 \sup_{|t|\le x^{C'}}
 \left|\sum_{d\in I}\mu(d)R_Z(d)\chi(d)d^{it}\right|
 \ll_{A',B',C',\delta}D_1(\log x)^{-A'}.
\end{equation}
For a bounded-variation weight $w$, (2.11), p.~6, multiplies this
bound by $\|w\|_\infty+\operatorname{Var}_I(w)$.
Fixed logarithmic powers in these norms are absorbed by a larger
fixed precision.
\end{inputthm}

\subsection{The energy theorem}\label{subsec:A}
\begin{theorem}[Theorem A]\label{thm:A}
Fix $A$ and $A_1\ge1$, and put
\begin{equation}\label{eq:A-constants}
 \begin{aligned}
 D&=\left\lceil\frac32(A+B_1+B+20)\right\rceil,
 & C_3&=\max(B_1,2B)+1,\\
 C_{\mathrm{zf},2}&=\max(D+4,C_3+1),
 & C_{\mathrm{zf},1}&=12(D+4),\\
 A_s&=A+B_1+B+10.
 \end{aligned}
\end{equation}
In the setting of Section~\ref{subsec:arc-reduction},
$\mathrm H1^\star(C_{\mathrm{zf},1},C_{\mathrm{zf},2};A_1)$ at $\GN$
implies \eqref{eq:alpha-star}, for a threshold depending on all
fixed parameters and $A_1$.
\end{theorem}
\begin{proof}
Split the characters into
$\mathcal X_s=\{\chi:\cond(\chi\psi)\le L^{A_1}\}$ and its
complement $\mathcal X_b$.
For $\chi\in\mathcal X_s$, let $\xi^*$ induce $\xi=\chi\psi$.
On the coefficient support,
\begin{equation}\label{eq:small-conductor-M}
 M_\xi(t)=\frac1M\sum_m\mu(m)R_Z(m)\ind_{(m,\GN)=1}
                    \xi^*(m)m^{i(t+v)}.
\end{equation}
Removing the unit indicator changes this by at most
\[
 \sum_{\substack{p\mid\GN\\p>Z}}\frac{2M/p}{M}
 \le2\omega_{\rm all}(\GN)/Z\le2Le^{-L^{0.24}}.
\]
Apply LWW Proposition~2.2, (2.5), p.~5, and (2.11), p.~6, to the
sum without the indicator: the modulus is at most $L^{A_1}$,
$M\ge x^\delta$ and $|t+v|\le x$ for large $x$.
For every fixed $A'$ this gives
$\sup_{|t|\le T_0}|M_\xi(t)|\ll L^{-A'}$.
This is the use of the sifted M\"obius input, whose proof cites
\cite[Lemma~3.9(ii)]{MSTT}.
For fixed $\psi$, the primitive inducing character determines
$\chi$, so $|\mathcal X_s|\le L^{2A_1}$.
The polynomial $B_{\xi_b}N_\xi$ is supported on $u\le32YN$ with
coefficients at most $C_BL\tau(u)/(16YN)$ and squared coefficient
norm $\ll L^5/(YN)$. The ordinary mean value theorem gives
$\int_{|t|\le T_0}|BN|^2\ll(T_0+YN)L^5/(YN)\ll L^5$.
Thus with $A'=A+A_1+4$,
\begin{equation}\label{eq:small-conductor-energy}
 \frac1{\ph(\GN)}\sum_{\chi\in\mathcal X_s}\int|BMN|^2
 \ll\frac{L^{2A_1-2A'+5}}{\ph(\GN)}
 \ll L^{-A}/\GN,
\end{equation}
using $\GN/\ph(\GN)\ll\log L$.

For $\chi\in\mathcal X_b$,
$q^*\mid k\GN$ and $L^{A_1}<q^*\le k\GN$.
The assumed region and Corollary~\ref{cor:P} yield
$|N_\xi(t)|\le9L^{1-D}$ for $|t|\le T_0+1$.
Discretize on this subset as above; restricting characters only
reduces the bounds for $\mathcal P$ and $E$.
The bulk is at most
\[
 L^{-2A_s+B_1+B+8}\ph(\GN)/\GN
 \le L^{-A-1}\ph(\GN)/\GN.
\]
If $M\ge\GN^{3/8+\eta_0}$, \eqref{eq:EM} gives
\begin{equation}\label{eq:caseA}
 \sum_E|BMN|^2
 \le C_B^2\sup_E|N|^2\sum_E|M|^2
 \ll L^{2-2D}\log L.
\end{equation}
If $M<\GN^{3/8+\eta_0}$, then
\begin{equation}\label{eq:caseB-length}
 \begin{gathered}
 N=\frac{MN}{M}\gg\frac{x}{M}
       \gg\frac{\GN^{5/8-\eta_0}}{L^B},\\
 N\ge\GN^{1/2+2\eta_0}\quad\text{for sufficiently large }\GN.
 \end{gathered}
\end{equation}
H\"older with exponents $3,3/2$, followed by
\eqref{eq:EM3} and \eqref{eq:EN}, gives
\begin{equation}\label{eq:caseB}
 \begin{split}
 \sum_E|B|^2|M|^2|N|^2
 &\le C_B^2\sup_E|N|^{2/3}\sum_E|M|^2|N|^{4/3}\\
 &\le C_B^2(9L^{1-D})^{2/3}
       (\sum_E|M|^6)^{1/3}(\sum_E|N|^2)^{2/3}
 \ll L^{(2/3)(1-D)+3+2}.
 \end{split}
\end{equation}
In either case the exceptional contribution is
\begin{equation}\label{eq:exceptional-end}
 \ll L^{6-2D/3}\le L^{-A-B_1-B-14}\le L^{-A-1}.
\end{equation}
One may also use the looser exponent
$(2/3)(1-D)+6$ in the second case and
$3-2D$ with a factor $\log L$ in the initial case; both leave the
same required margin. Sum the two families and divide by $\ph(\GN)$:
\[
 \frac1{\ph(\GN)}\sum_{\chi\in\mathcal X_b}\int|BMN|^2
 \ll L^{-A-1}/\GN+L^{-A-1}/\ph(\GN)\ll L^{-A}/\GN.
\]
\end{proof}
The inputs to Theorem~\ref{thm:A} are the ordinary mean value theorem,
P\'olya--Vinogradov, Burgess at $r=2$, partial fractions and zero
counting, Perron, and LWW Proposition~2.2. The half-plane
\eqref{eq:H1} enters only by supplying \eqref{eq:Hstar}.
For one fixed region the theorem supplies its selected fixed
precision $A$; it does not assert all precisions from one fixed width.

\subsection{The small-radical input}\label{subsec:radical}
\begin{inputthm}[Iwaniec, 1974, Theorem 2]\label{input:I}
For $q\ge3$, put $d=\rad(q)$,
$\ell=\log(q(|t|+3))$ and
\begin{equation}\label{eq:Iwaniec}
 \vartheta^{-1}=4\cdot10^4
       \{\log d+(\ell\log(2\ell))^{3/4}\}.
\end{equation}
There is at most one nonprincipal character modulo $q$ and one
zero $\rho_0$ in $\Re\rho_0>1-\vartheta$; if present, the
character is real and the zero real and simple
\cite[Theorem~2, p.~97]{Iwaniec1974}.
\end{inputthm}
\begin{lemma}[Lemma Q]\label{lem:Q}
If $\rad(\GN)\le L^C$ and $k\le L^{A_0}$, every real primitive
character of conductor $q^*\mid k\GN$ has
\begin{equation}\label{eq:Q}
 q^*\le8\rad(k\GN)\le8L^{A_0+C}\le L^{A_1},
 \qquad A_1=A_0+C+1,\quad L\ge8.
\end{equation}
\end{lemma}
\begin{proof}
A real primitive character is a Kronecker symbol with fundamental
discriminant of absolute value $m$, $4m$ or $8m$, with $m$ odd and
squarefree \cite[Chapter~5]{Davenport}. Hence
$q^*\mid8\rad(q^*)\mid8\rad(k\GN)$.
\end{proof}
\begin{corollary}[Small-radical region]\label{cor:small-radical}
For $\rad(\GN)\le L^C$ and $A_1=A_0+C+1$,
\eqref{eq:Hstar} holds for every fixed
$C_{\mathrm{zf},1},C_{\mathrm{zf},2}$ once
$L\ge L_0(C_{\mathrm{zf},1},C_{\mathrm{zf},2},C,A_0)$.
\end{corollary}
\begin{proof}
For $q^*\mid k\GN$, $q^*>L^{A_1}$, apply
Input~\ref{input:I} with $q=q^*$.
Here $\rad(q^*)\le L^{A_0+C}$ and
$\ell=\log(q^*(|t|+3))\le2L$ at the fixed height $L^{C_{\mathrm{zf},2}}$.
Thus
\begin{equation}\label{eq:Iwaniec-width}
 \begin{split}
 \vartheta^{-1}&\le4\cdot10^4
 \{(A_0+C)\log L+(2L\log(4L))^{3/4}\},\\
 \vartheta&\ge c_I L^{-3/4}(\log L)^{-3/4}
       > C_{\mathrm{zf},1}\log L/L
 \end{split}
\end{equation}
for sufficiently large $L$, with absolute $c_I>0$ of size $10^{-5}$.
A real exceptional character would contradict Lemma~\ref{lem:Q},
so the exceptional case is absent. The strict width comparison
also gives the closed boundary in \eqref{eq:Hstar}.
\end{proof}
The explicit radical dependence is needed because $k$ contributes
radical up to $L^{A_0}$ even when $\GN=2^m$; a prime-power region
alone does not cover these conductors.

\begin{corollary}[The analytic alternatives]\label{cor:alternatives}
Theorem~\ref{thm:A} supplies \eqref{eq:alpha-star} from the selected
$\mathrm H1^\star$, and therefore from $\mathrm H1$ for every even
$\GN$. In the class $\rad(\GN)\le L^C$ it follows from
Corollary~\ref{cor:small-radical}. For $0<|h|\le L^C$, all
characters have conductor at most $L^{A_0+C}$, but the scale-$1$
Mellin reduction does not apply; Section~\ref{subsec:small-major}
uses LWW Section~5 instead.
\end{corollary}

\subsection{The classical-region obstruction}\label{subsec:obstruction}
\begin{remark}[A separate contour estimate]\label{rem:other-contour}
Fix $C_t\ge2$. Under \eqref{eq:H1}, formula~(5.15) of LWW,
p.~44, separates a prime in $\mathcal P_1$ from $B_{\xi_b}$.
Lemma~3.5 of \cite{GuoRGP}, followed by dyadic summation and
partial summation, gives
\begin{equation}\label{eq:other-prime-estimate}
 \sum_{p\in\mathcal P_1}\frac{\xi_b(p)p^{it'}}p
 \ll L^{C_{\mathrm{prime}}+1}\exp(-L^{a_1}/16)
\end{equation}
uniformly for $|t'|\le\GN^{C_t}$ and nonprincipal $\xi_b$ of
modulus $k\GN\le\GN^2$. Here $C_{\mathrm{prime}}$ is fixed,
depending on $C_t$ and the fixed exponent ranges.
The cited Perron argument uses $\Re s=15/16$ and
$L'/L\ll\log^2(q(|t|+3))$.

For the classical real-zero lower bound
$1-\beta_1\gg_\varepsilon q^{-\varepsilon}$, even at
$q\le\exp(L^{a_1/2})$ the guaranteed lower bound for
$(1-\beta_1)L^{a_1}$ can tend to zero. It gives no divergence to
infinity, and leaves a possible exceptional conductor in
$(L^{C},\exp(L^{a_1/2})]$ outside this argument.
\end{remark}

\begin{proposition}[Obstruction (EC)]\label{prop:EC}
Let $\chi\bmod q$ be real primitive, with $q\mid\GN$.
For $p\nmid\GN$ one has
\begin{equation}\label{eq:EC-identity}
 \mu(\GN-p)=\chi(-1)\chi(p)\varepsilon(\GN-p)\mu^2(\GN-p),
 \qquad \varepsilon(n)=\prod_{p'\mid n}(-\chi(p')).
\end{equation}
Suppose every split prime ($\chi(p')=1$) at most $\log q$ divides
$\GN$, and at most $\pi(\GN)/20$ primes $p<\GN$ have either
$\chi(p)=1$ or a split prime factor exceeding $\log q$ in $\GN-p$.
Then, for sufficiently large such $\GN$,
\begin{equation}\label{eq:EC}
 \left|\sum_{p<\GN}\mu(\GN-p)\right|\ge0.27\pi(\GN).
\end{equation}
\end{proposition}
\begin{proof}
For squarefree $n$ coprime to $q$,
$\varepsilon(n)=\mu(n)\chi(n)$; for nonsquarefree $n$ both sides of
\eqref{eq:EC-identity} vanish after multiplication by $\mu^2(n)$.
Also $\chi(\GN-p)=\chi(-1)\chi(p)$ when $p\nmid\GN$,
proving the identity.
Outside the stated exceptional set, $\chi(p)=-1$.
No split prime can divide $\GN-p$: a small split prime divides
$\GN$ and would then divide $p$, contrary to $p\nmid\GN$;
a larger split prime is excluded by hypothesis.
Thus $\varepsilon(\GN-p)=1$, and each squarefree such value has
$\mu(\GN-p)=-\chi(-1)$.
Section~\ref{sec:density}, whose proof uses only Siegel--Walfisz,
gives $(A_{\GN}+o(1))\pi(\GN)$ squarefree values uniformly.
Removing the exceptional set loses at most $\pi(\GN)/20$ of them,
and their possible opposite contribution loses another
$\pi(\GN)/20$. The $O(\omega_{\rm all}(\GN))$ nonunit primes are
negligible. Hence the left side of \eqref{eq:EC} is at least
$(0.3739-2/20+o(1))\pi(\GN)>0.27\pi(\GN)$.
\end{proof}
This proposition establishes an obstruction under its split-prime
hypotheses. The classical lower bound
$1-\beta\gg_\varepsilon q^{-\varepsilon}$ does not by itself
exclude that scenario for a large conductor dividing $\GN$.
The possible zero-size condition $(1-\beta)\log\GN\to0$
does not replace the separate split-prime hypotheses in this proposition.
It identifies why the present uniform argument uses
\eqref{eq:Hstar}. In the small-radical class all the real
conductors dividing $k\GN$ are small by Lemma~\ref{lem:Q}.

\section{Type I and the common divisor}\label{sec:typeI}
\begin{proposition}[Type I for the Goldbach shift]\label{prop:I}
For the bounded function in Proposition~\ref{prop:vaughan},
LWW Proposition~3.1 with the following changes gives
$T_1(F),T_2(F)\ll_A xL^{-A}$, uniformly in the range
\eqref{eq:ranges}.
\end{proposition}
\begin{proof}
Use the argument of LWW Section~3, pp.~10--14.
The classes with shift $+2$ become, for example,
$k\equiv\GN\bar q\pmod d$; every divisor $d$ built from list
primes is coprime to $\GN$.
The sifted modulus $q$ can have a prime $>Z$ in common with $\GN$.
Discarding those terms costs
\begin{equation}\label{eq:I-q-discard}
 \ll xL^3\omega_{\rm all}(\GN)/Z=O_A(xL^{-A}).
\end{equation}
At this point the unit condition on the sifted inner variable can
also be removed at that cost. It is removed before expanding
$R_Z$, so that no additional condition on the later M\"obius variable
is imposed.

After the truncated inclusion--exclusion in LWW Section~3, p.~11,
$r$ is odd and squarefree with small prime factors, and may share
primes with $\GN$. Put $g=(r,\GN)$.
The list divisor $d$ is coprime to $r\GN$, so $dg$ is squarefree.
Use the exact identity (3.4), p.~12, in the form
\begin{equation}\label{eq:mu-dg}
 \mu(dg v)=\mu(dg)
     \sum_{\substack{a\mid v\\p\mid a\Rightarrow p\mid dg}}\mu(v/a).
\end{equation}
The substitution is now
\begin{equation}\label{eq:I-substitution}
 \GN-qrj=dgat,\qquad
 t\equiv(\GN/g)(da)^{-1}\pmod{qr/g},\qquad t\ge1.
\end{equation}
The inverse exists: primes of $a$ lie in $dg$, so
$(da,qr/g)=1$. Indeed $d$ is disjoint from $qr$ and the primes of
$g$ have been removed from the squarefree $r$; $q$ is $Z$-rough.
Also $(\GN/g,r/g)=1$, $(\GN,q)=1$ and $(g,qr/g)=1$.
Thus the class is reduced. Primes of $a$ may divide $g$; no
condition $(g,da)=1$ is needed or imposed.
Since $qr,dga$ are odd and $\GN$ is even,
$j$ odd is equivalent to $t$ odd.
The M\"obius variable therefore lies in a reduced class modulo
$2qr/g$, with $0$ excluded.

Apply Bombieri--Vinogradov in LWW (3.1), p.~10, with the
$\tau$-weighted grouping of p.~13.
For each fixed $g$, $qr/g$ determines $qr$ and its rough part $q$
and small-prime part $r$ uniquely; one may bound coefficients by
$L$ or by $L\tau(qr/g)$ without a $\tau(g)$ loss.
The level is comparable to $\GN/(dga)$ and depends on $g$.
Apply the bound separately for each $g$ and then sum its reciprocal
weight:
\begin{equation}\label{eq:I-g-sum}
 \sum_{g\mid\rad(\GN)}\frac1g
 =\prod_{p\mid\GN}(1+1/p)\ll\log L.
\end{equation}
The sum over $a$ of $1/a$ has the same admissible Euler-product
bound as in LWW. For the large-$a$ part in (3.5), p.~13, the new
Rankin factor is
\begin{equation}\label{eq:I-Rankin}
 \prod_{p\mid g}(1-p^{-1/20})^{-1}
 \le19^{R''}=L^{O(1)},\qquad R''=\lceil C'\log L\rceil,
\end{equation}
absorbed by $x^{-1/50}$. Here $r$ has at most $R''$ prime factors.
The incompatibility $(r,d)>1$ is unchanged: a list prime in both
would divide $\GN$.

The positive cofactor interval is
\begin{equation}\label{eq:I-endpoints}
 \left[\frac{\GN-2x}{dga},\frac{\GN-x}{dga}\right),
 \qquad \GN/(dga)\ll L^B x/(dga).
\end{equation}
The maximum over endpoints in Bombieri--Vinogradov handles this
interval. Replacing $x$ by $\GN$ in the level costs only the fixed
power $L^B$, covered by the chosen precision.
Partial summation for $T_1$ is that of LWW Section~3, p.~14.
\end{proof}

\section{Removal of the auxiliary weights}\label{sec:weights}
\begin{proposition}[Normalized weight removal]\label{prop:weights}
For each fixed $b$ choose $b$ disjoint copies $W_j$ of
\eqref{eq:weight}, and let $m_j$ be their reduced-class means for
$W_j(\GN-a)$. Then
$m_j=2^{K(\rho-1)}+o(1)>0$, and
\begin{equation}\label{eq:weight-removal}
 \begin{split}
 \limsup_{\GN\to\infty}\frac1{\GN}
       \left|\sum_{p<\GN}\mu(\GN-p)\log p\right|
 &\le\sqrt{C_{K,\rho}/b}\\
 &\quad+\limsup_{\GN\to\infty}
 \max_{1\le j\le b}\frac{
     |\sum_{p<\GN}\mu(\GN-p)W_j(\GN-p)\log p|}{m_j\GN}.
 \end{split}
\end{equation}
The constant $C_{K,\rho}$ is independent of $b$.
\end{proposition}
\begin{proof}
Let $P_0$ be the product of the list primes for a fixed copy.
For every $d\mid P_0$,
\begin{equation}\label{eq:CRT-weight}
 \frac1{\ph(P_0)}\sum_{\substack{a\bmod P_0\\(a,P_0)=1}}
          \ind_{d\mid\GN-a}=\frac1{\ph(d)}.
\end{equation}
This follows because $\GN$ is a unit modulo $d$, so the divisibility
condition selects one reduced class. It does not require translation
to permute all reduced classes.
Thus the mean (2.14) and second moment (2.15) of LWW, p.~7, apply.
The truncated divisor expansion and Bombieri--Vinogradov for
$p\equiv\GN\pmod d$, $(d,\GN)=1$, give the same weighted
moments over primes. All moduli are odd squarefree products of
list primes, and the maximum over reduced classes makes the estimate
uniform in $\GN$.

For the weight truncation write $R_w=\lceil\log L\rceil$, as in
LWW Section~2.3, p.~9. Its integer-average error in this range is
\begin{equation}\label{eq:weight-truncation}
 O(L^{B+1}/R_w!)=o(1),\qquad
 R_w!=L^{(1+o(1))\log\log L}.
\end{equation}
This distinguishes $R_w$ from the moment order $R$.
The variance argument (2.16)--(2.19), pp.~7--10, is applied to
$b^{-1}\sum_jW_j/m_j$. Cauchy--Schwarz against the bounded target
and the triangle inequality give \eqref{eq:weight-removal}.
The means tend to the same positive fixed constant, so normalization
has a bounded cost independent of $b$.
For each fixed $b$, let $\GN$ tend to infinity before letting $b$
tend to infinity; thresholds may depend on its chosen bands.
\end{proof}

\section{Squarefree values}\label{sec:density}
\begin{proposition}[Uniform squarefree density]\label{prop:squarefree}
For even $\GN\to\infty$,
\begin{equation}\label{eq:squarefree-weighted}
 \sum_{\substack{p<\GN\\p\nmid\GN}}\mu^2(\GN-p)\log p
 =A_{\GN}\GN+O_A(\GN L^{-A}),
 \qquad A_{\GN}\ge0.3739.
\end{equation}
Consequently
\begin{equation}\label{eq:squarefree-count}
 \#\{p<\GN:\mu^2(\GN-p)=1\}=(A_{\GN}+o(1))\pi(\GN),
\end{equation}
uniformly in even $\GN$.
\end{proposition}
\begin{proof}
Use $\mu^2(n)=\sum_{d^2\mid n}\mu(d)$.
If $p\nmid\GN$ and $d^2\mid\GN-p$, then $(d,\GN)=1$.
For $d\le L^{A+2}$, Siegel--Walfisz in the reduced class
$p\equiv\GN\pmod{d^2}$, at a sufficiently large fixed precision,
gives main term
\begin{equation}\label{eq:squarefree-series}
 \GN\sum_{(d,\GN)=1}\frac{\mu(d)}{\ph(d^2)}=A_{\GN}\GN,
\end{equation}
with extension error $O(\GN L^{-A-1})$.
The tail has the finite range required by $\GN-p>0$:
\begin{equation}\label{eq:squarefree-tail}
 \sum_{L^{A+2}<d\le\sqrt{\GN}}(\GN/d^2+1)L
 \ll\GN L^{-A-1}+\sqrt{\GN}L.
\end{equation}
The last term is absorbed in $O_A(\GN L^{-A})$.
This is LWW Proposition~6.1 and (6.3), pp.~46--47, with the size
and reduced-class checks for the varying shift supplied here.
Primes dividing $\GN$ cost $O(\omega_{\rm all}(\GN)L)=O(L^2)$.
To pass to ordinary counts, use Lemma~\ref{lem:unweight} below
with the bounded target $\mu^2(\GN-p)$.
The same progression proof permits any endpoint $y\le\GN$.
\end{proof}

\begin{lemma}[Removing the prime logarithm]\label{lem:unweight}
For any bounded target $|f(p)|\le1$ on primes below $\GN$, possibly
depending on $\GN$, and fixed $B>2$,
\begin{equation}\label{eq:unweight}
 \sum_{p<\GN}f(p)
 =\frac1L\sum_{p<\GN}f(p)\log p
 +O_B\left(\frac{\GN\log L}{L^2}+\frac{\GN}{L^B}\right).
\end{equation}
\end{lemma}
\begin{proof}
For $p>\GN/L^B$, $\log p/L=1+O_B(\log L/L)$.
Chebyshev's bound gives $O(\GN/L)$ such primes.
For the remaining primes bound both sums by their number, at most
$\GN/L^B$; this proves the identity without varying the target
sequence in a partial-summation limit.
\end{proof}

\section{Assembly of the Goldbach statements}\label{sec:assembly}
\begin{proof}[Proofs of Theorems~\ref{thm:G} and~\ref{thm:S}]
Choose $\rho$, $A^*>0$ and then $B>A^*+3$.
Let $C_{\mathrm{coef}}$ be the corresponding fixed coefficient and
smoothing-loss exponent for this choice of $B$.
Complete the parameter choices in the following order:
\begin{equation}\label{eq:parameter-order}
 \begin{gathered}
 \rho\ ;\ A^*\ ;\ B\ ;\ C_{\mathrm{coef}}\ ;\\
 D^*>2(A^*+B+C_{\mathrm{coef}}+c_*)\ ;\ E_0\ ;\ G\ ;\\
 A_0\ ;\ K\ ;\ C^*,C^\dagger,C_0\ ;\ B_0,B_1\ ;\\
 A_{\mathrm{energy}}=D^*+O(A_0+B_1)\ ;\\
 D,C_{\mathrm{zf},1},C_{\mathrm{zf},2},C_3,A_s,A_1\ ;\\
 \text{the fixed precision of Proposition 2.2}\ ;\ a_i\ ;\ L_0.
 \end{gathered}
\end{equation}
Here $c_*$ is a fixed reserve for scale and dyadic losses, chosen
large enough before the later parameters.
The choices for the moment estimate are those in LWW
Section~4.7, p.~39; $A_0,K$ are independent of the band locations.
The zero-free constants are subsequently defined by
\eqref{eq:A-constants} with $A=A_{\mathrm{energy}}$ and do not
enter $C_{\mathrm{coef}}$ or the inequality for $D^*$.
For any fixed $b$ the $bK$ band exponents can be selected after these
saving exponents; their positions affect thresholds, not these choices.

On each of the $O(\log L)$ dyadic prime intervals, the Vaughan
reduction and Proposition~\ref{prop:I} reduce the weighted prime sum
to Type~II sums $T_Y$ on $O(L^{0.2})$ product scales.
Propositions~\ref{prop:II} and~\ref{prop:arc-reduction}, followed by
Theorem~\ref{thm:A}, give
\begin{equation}\label{eq:assembly-square}
 \QQ=\Qmaj+O(\GN YL^{-D^*}),\quad
 |\Qmaj|\ll\GN YL^{-D^*},\quad
 |T_Y|^2\ll(\GN/Y)\QQ\ll\GN^2L^{-D^*}.
\end{equation}
Thus $|T_Y|\ll\GN L^{-D^*/2}$, and the larger fixed budget in
\eqref{eq:parameter-order}, using $\GN/X\ll2L^B$, gives
$T_3(F)\ll xL^{-A^*}$.
For Theorem~\ref{thm:G}, use the selected $\mathrm H1^\star$;
\eqref{eq:H1} supplies it for every fixed selection of constants.
For Theorem~\ref{thm:S}, use Corollary~\ref{cor:small-radical} with
$A_1=A_0+C+1$ and the Burgess input in Theorem~\ref{thm:A}.
In either case
\begin{equation}\label{eq:weighted-conclusion}
 \sum_{x<p\le2x}(\log p)\mu(\GN-p)W(\GN-p)=O(xL^{-A^*}),
 \qquad
 \sum_{p<\GN}(\log p)\mu(\GN-p)W(\GN-p)=o(\GN).
\end{equation}
The omitted small primes have the negligible bounded-weight cost
stated in Remark~\ref{rem:ranges}.

For each fixed $b$, every normalized weighted limit in
\eqref{eq:weight-removal} is zero.
Then let $b\to\infty$ to get
$\sum_{p<\GN}\mu(\GN-p)\log p=o(\GN)$.
Lemma~\ref{lem:unweight} and $\pi(\GN)\sim\GN/L$ give the initial
conclusion of \eqref{eq:G}.
Finally, for $\kappa=\pm1$,
\begin{equation}\label{eq:mu-sign-identity}
 \#\{p<\GN:\mu(\GN-p)=\kappa\}
 =\tfrac12\left(\#\{p<\GN:\mu^2(\GN-p)=1\}
                 +\kappa\sum_{p<\GN}\mu(\GN-p)\right).
\end{equation}
Proposition~\ref{prop:squarefree} proves the sign densities.

All comparisons use $\log x=L+O_B(\log L)$ on the stated range.
Puncturing, the large common-prime discard, the power errors in
Lemma~\ref{lem:R} and the $O(\log L)$ dyadic summation are uniform
in even $\GN$. Constants depend on the fixed bands, $K,\rho,b$,
and requested precisions, and in the small-radical case on $C$.
The limits are ineffective through Siegel--Walfisz and
Proposition~2.2. A fixed $\mathrm H1^\star$ supplies the selected
fixed energy precision; the half-plane and the small-radical region
supply every fixed precision needed in this order.
\end{proof}

\section{Small shifts}\label{sec:small}
Throughout this section $x$ is the large parameter, $L=\log x$,
and $0<|h|\le L^C$, with fixed $C$. All constants are uniform in
this range. We use the prime form $p+h$ on the domain $p+h\ge1$.

\subsection{Setting and reductions}\label{subsec:small-reduction}
Define $F_h(n)=\mu(n+h)W(n+h)$ when $n+h\ge1$, and $0$ otherwise.
On $(x,2x]$ the arguments are positive for large $x$.
Puncture the bands at $p\mid h$, costing at most
$\omega_{\rm all}(h)e^{-L^{0.1}}$ in reciprocal mass.
Every prime divisor of $h$ is below $Z$, so the sifted Vaughan
variables satisfy $(mnrs,h)=1$ automatically.
Write $mn+h=ah'$, $a\in\Aset$. A common prime of $a,m$ would
divide $h$, which puncturing excludes.
Remove $(a,h')>1$ as in LWW (4.3), p.~15.
Now $mn+h\asymp X$ on every rectangle, so the common cofactor
range is $O(X/Y)$ and
\begin{equation}\label{eq:small-cauchy}
 |T_Y|^2\ll(X/Y)Q_Y,
 \qquad Q_Y\ge0.
\end{equation}
The Lemma~4.1 check is
$d\mid\ell h'+h\Rightarrow(d,\ell)\mid h$ and
$(\ell,h)=1$, so $(d,\ell)=1$.
For $(a,b)=1$ the equation is
\begin{equation}\label{eq:small-det}
 b(mn+h)=a(rs+h)
 \quad\Longleftrightarrow\quad
 t=bmn-ars=-h(b-a),\qquad h\mid t,\quad t''=-t/h=b-a.
\end{equation}
It is LWW (4.8), p.~16, when $h=2$.

\subsection{Oriented bases, parity and root counts}\label{subsec:small-basis}
\begin{lemma}[Small-shift blocks and orientation]\label{lem:small-basis}
Use labelled lattices $\Lambda_\varrho$ of index $|h|$ for unit
slopes modulo $|h|$ and odd-odd endpoint positions.
For a primitive odd-odd root $P=(u,v)^T$ with $(uv,h)=1$, choose
$ud-vc=1$, $0\le c<u$, $r_P=c/u$. Put
\begin{equation}\label{eq:small-basis}
 \begin{split}
 P^*&=-h(c,d)^T,\qquad
 g_h(P)=\begin{pmatrix}u&-hc\\v&-hd\end{pmatrix}
       =\begin{pmatrix}u&-hur_P\\v&-h(1/u+vr_P)\end{pmatrix},\\
 E_h&=\operatorname{diag}(1,-h),\quad
 \det g_h(P)=-h,\quad \tau_h(z)=z_1-hr_Pz_2.
 \end{split}
\end{equation}
This is a basis of the index-$|h|$ lattice. For $Q=g_h(P)z$,
\begin{equation}\label{eq:small-coordinates}
 \det(P,Q)=-hz_2,\qquad Q_1=u\tau_h(z),\qquad
 Q_2=v\tau_h(z)-hz_2/u.
\end{equation}
Lattice-primitivity is $(z_1,z_2)=1$; integer-plane primitivity
also requires $(Q_1,h)=1$.
For primitive $z$ and $p\nmid h$, list divisibility is
$[z]_p=\ker(u,-hc)\bmod p$.
For $\det(z,w)=1$, the intrinsic ratio obtained with a complement
of determinant $|h|$ is
\begin{equation}\label{eq:small-ratio}
 \beta_Q^{(h)}\equiv-\operatorname{sgn}(h)
      \frac{\tau_h(w)}{\tau_h(z)}\pmod1,
 \qquad
 \partial_r\frac{\tau_h(w)}{\tau_h(z)}
        =-\frac{h}{\tau_h(z)^2}.
\end{equation}
Its absolute derivative lies in $[|h|/256,256|h|]$.
The odd-odd restriction is
\begin{equation}\label{eq:small-parity}
 \begin{cases}
 z_1\text{ odd},&h\text{ even},\\
 z_1\text{ odd and }z_2\text{ even},&h\text{ odd}.
 \end{cases}
\end{equation}
\end{lemma}
\begin{proof}
The columns lie in $\Lambda_\varrho$ and have determinant $-h$,
so the basis and coordinate assertions follow with that orientation.
The first-row kernel assertion follows from invertibility modulo
$p\nmid h$. If $W=g_h(P)w$ has relative determinant $1$,
then $\det(Q,W)=-h$. Multiplication of $W$ by
$-\operatorname{sgn}(h)$ gives determinant $|h|$, proving
\eqref{eq:small-ratio}. At an integer-plane primitive state this
ratio is $|h|r_Q\pmod1$.
Both good-state tests use distance to integers, so changing the
ratio's sign preserves them; absolute-value boundary bounds also
preserve their estimates.
The derivative formula follows by differentiation, and
$\tau_h(z)\in[1/16,16]$ gives its range, bounded below by $1/256$
and above by $x^{o(1)}$.
For even $h$, $P^*\equiv0\pmod2$, giving the initial parity case.
For odd $h$, $P^*\equiv(c,d)\pmod2$ and
$d-c\equiv ud-vc\equiv1\pmod2$. Hence
$Q\equiv(z_1+cz_2,z_1+dz_2)\pmod2$, proving the other case.
\end{proof}

Let $\mathfrak M_h=\{\theta:|h|\theta\in\Mone\}$.
The covering identity is
\begin{equation}\label{eq:small-cover}
 \int_{\mathfrak M_h}e(\theta(t+h(b-a)))\,d\theta
 =\ind_{h\mid t}\int_{\Mone}e(\phi(t''-b+a))\,d\phi.
\end{equation}
Indeed the branch substitution first gives phase
$e(-\operatorname{sgn}(h)\phi(t''-b+a))$; for $h>0$ reflect
$\phi$, using symmetry of $\Mone$. The kernel is block diagonal
and its reduced determinant is exactly $\det(z,z')/D$.
A full block can contain $\asymp d_0YMN/|h|$ positions, so
Lemma~\ref{lem:block-size}'s $O(d_0Y)$ application bound is replaced
by the local partition below.

Parity is imposed only through the multiplier $\chi$ in LWW
(4.34), p.~25, and the corresponding sublattice restriction in
$M_R$: use \eqref{eq:small-parity} on the full relative integer
lattice. This projection has norm at most $1$.
The labelled enlargement is used only for nonunit states modulo
$h$, where the endpoint vector vanishes. In particular even positions
are not separately reintroduced this way. For odd $h$, allowed
reduced determinants are even, which is just a restriction of this
projection. The endpoint multiplier $\chi(B_Dk)$ in (4.45), p.~29,
and compression of the clean-edge operator cannot increase its norm.

\begin{lemma}[Lemma R for small shifts]\label{lem:Rh}
Let $d_0YM,N\ge x^\delta$ and let $S$ be odd squarefree,
$(S,h)=1$, $S\le\exp(C_0L^{0.98})$.
Use the intervals $I_1,I_2,I_3$ of Lemma~\ref{lem:R}.
For a matrix $g_0\bmod S$ of determinant $-h$, the number of
primitive roots with $u,v$ odd, $(uv,h)=1$,
$u/(d_0YM)\in I_1$, $v/N\in I_2$, $r_P\in I_3$ and
$g_h(P)\equiv g_0\pmod S$ is
\begin{equation}\label{eq:Rh-count}
 \frac{c_h d_0YMN|I_1||I_2||I_3|}
      {|\operatorname{SL}_2(\Z/S\Z)|}
       +O(d_0YMNx^{-c\delta}),\qquad
 c_h=\frac1{\zeta(2)}\prod_{p\mid2h}\frac{p-1}{p+1}.
\end{equation}
Prescribed lines have weights $1/(p+1)$ and a uniform error of the
same form, with a possibly smaller fixed positive $c$.
For $h=\pm2$, $c_h=1/(3\zeta(2))$.
\label{eq:ch}
\end{lemma}
\begin{proof}
Use LWW Lemma~A.1 and (A.5), pp.~47--51, with the original ratio
$r_P=c/u$. Since $(h,S)=1$,
$g_h\equiv g_0\pmod S$ is equivalent to
$g\equiv g_0E_h^{-1}\pmod S$.
The odd-odd condition supplies the original parity factor at $2$.
Sieve the odd part of $\rad(h)$ on both coordinates as in
Lemma~\ref{lem:R}; its local factor is $(p-1)/(p+1)$ and its
divisor cost $2^{\omega_{\rm all}(h)}\le L^{O(1)}$.
The swapped-coordinate case is LWW's determinant-one calculation
with displacement $1/(uv)$, since the counted ratio remains $r_P$.
Thus there is no $\GN/(uv)$ displacement to discard.
The ratio map $r\mapsto|h|r$ preserves Haar measure, and the
line distribution follows from the determinant-$(-h)$ matrix
coset just as in \eqref{eq:R-lines}.
\end{proof}

\subsection{Locality within a small-shift block}\label{subsec:small-locality}
\begin{proposition}[Small-shift Type II estimate]\label{prop:small-II}
Using $\mathrm H2'$ and the changes below,
\begin{equation}\label{eq:small-II}
 |Q_Y-Q_Y^{\mathrm{maj},h}|\ll_A XYL^{-A}
\end{equation}
uniformly for $0<|h|\le L^C$.
\end{proposition}
\begin{proof}
The relative kernel has $t''=\det(z,z')/D$ and the same list
structure. Rows~5, 9 and 17 of Table~\ref{tab:dictionary} are
replaced as follows; all support-dependent constants refer to the
fixed cutoff $\psi_1$.
On an edge,
\begin{equation}\label{eq:small-slopes}
 \begin{split}
 |\det(P_j,P_{j+1})|&\ll|h|d_0Y,\\
 |\operatorname{slope}(P_{j+1})-\operatorname{slope}(P_j)|
       &\ll\frac{|h|}{d_0YM^2}.
 \end{split}
\end{equation}
Summing at most $2R$ edges from $P=P_0$ gives
\begin{equation}\label{eq:small-path}
 |\det(P,P_j)|\ll R|h|d_0Y,\quad
 |(z_j)_2|\ll Rd_0Y,\quad
 |(z_j)_1|=|\tau_h(z_j)+hr_P(z_j)_2|
              \ll|h|Rd_0Y=x^{o(1)}.
\end{equation}
Use the relative cube $Z_0$ of radius $2C_{\rm pos}|h|Rd_0Y$ in
(4.33), p.~24. Its coefficient mass remains
$\exp(O_K(L^{0.73}))$, since
$\log(|h|Rd_0Y)=O(L^{0.21})$.
The second position boundary, with its orientation retained, is
\begin{equation}\label{eq:small-boundary}
 E_2=t_0\tau_h(z)-\frac{hz_2}{uN},\qquad
 \left|\frac{hz_2}{uN}\right|\ll\frac{R|h|}{MN}.
\end{equation}
Its derivatives have the bounds of LWW (4.30), pp.~23--24.

Inside each labelled block partition the physical slope $P_2/P_1$
into half-open intervals of length $|h|/(d_0YM^2)$.
By \eqref{eq:small-slopes}, each interval interacts with $O(1)$
neighbors in that block. Each contains $O(|h|d_0Y)$
lattice-primitive positions: choose one such $P$ in it.
Other positions have
$|\det(P,Q)|\ll|h|d_0Y$ and determinants in $|h|\Z$.
For a fixed determinant their differences belong to
$\Q P\cap\Lambda_\varrho=\Z P$, by lattice-primitivity,
so the first-coordinate strip gives $O(1)$ choices.
This reasoning also applies to added nonunit states and does not
assume their primitivity in $\Z^2$.
The trace factor is
\begin{equation}\label{eq:small-trace}
 (|h|d_0Y)^{1/(2R)}
       =\exp(O(L^{0.21}/L^{1/2}))=O(1).
\end{equation}
For the comparison count, $h\mid bk_0-a\ell_0$ and
$|bk_0-a\ell_0|\ll|h|Y$, giving $O(1)$ choices of $\ell_0$ in
one class modulo $|h|$ in an interval of length $O(|h|)$.
For endpoint primitivity, $p\mid(m,s)$ would give
$p\mid h(b-a)$ with $p>Z>|h|$ and
$0<|b-a|<4Y<Z$, a contradiction.
The other dictionary rows, Lemma~\ref{lem:Rh}, and LWW's arguments
then give \eqref{eq:small-II}.
\end{proof}

\subsection{Small-shift major arcs}\label{subsec:small-major}
\begin{proposition}[Small-shift major term]\label{prop:small-major}
Using the arguments of LWW Section~5 as written,
\begin{equation}\label{eq:small-major}
 |Q_Y^{\mathrm{maj},h}|\ll_D XYL^{-D}
\end{equation}
uniformly for $0<|h|\le L^C$.
\end{proposition}
\begin{proof}
The kernel from \eqref{eq:small-cover} carries
\[
 \ind_{bmn\equiv ars\ (h)}\psi_1(t''/Y)
       \int_{\Mone}e(\phi(t''-b+a))\,d\phi.
\]
Now $\psi_1(t''/Y)=\psi_1((ars-bmn)/(hY))$ forces
$|bmn-ars|\ll|h|Y\ll bmn$, a genuine localization.
Apply the proof of LWW Proposition~5.3, pp.~42--45, with these changes.
The congruence expands as
\begin{equation}\label{eq:small-character}
 \ind_{bmn\equiv ars\ (h)}
   =\frac1{\ph(|h|)}\sum_{\chi\bmod|h|}\chi(bmn)\overline{\chi(ars)}.
\end{equation}
On $\phi=h_1/k+\lambda_0/Y$,
\[
 e(h_1t''/k)=e(-h_1bmn/(hk))e(h_1ars/(hk)).
\]
Expand these unit-residue functions factor by factor in characters
modulo $|h|k$, and the factors $e(\mp h_1(b-a)/k)$ modulo $k$.
The total character mass is at most
$(|h|k)^6\le L^{6(A_0+C)}$.

Mellin separation uses $|h|Y$ in place of $Y$ and reflects the
cutoff when its orientation requires it. Its relative scale is
$|h|/X$, giving
\begin{equation}\label{eq:small-frequencies}
 |t|\le(X/|h|)L^{B_1}\le XL^{B_1},\qquad
 T_*=\exp(L/(\log L)^2).
\end{equation}
These are in the range of LWW Lemmas~5.1--5.2, pp.~40--41.
The kernel has $X$ replaced by $X/h$; its prefactor gains at most
an additional $|h|\le L^C$.
Fix one character term. Its conductor is at most
$|h|k\le L^{A_0+C}$.
LWW (5.14), p.~43, follows with arbitrary fixed precision $A$ from
(5.2), p.~40, Lemma~5.2 and Proposition~2.2 at the fixed modulus
exponent $A_0+C$.
Lemma~5.2 is uniform in the bounded character twist and gives an
energy estimate. The arbitrary saving comes from Proposition~2.2.
To absorb the character mass and the prefactor choose
\begin{equation}\label{eq:small-budget}
 A=D+9A_0+7C+B_1+1
\end{equation}
in place of LWW's $D+9A_0+B_1+1$.
The rest of their assembly is used as written.
All moduli are $L^{O(1)}$, so this step uses no
$\mathrm H1$ or $\mathrm H1^\star$.
It uses LWW Section~5 in addition to $\mathrm H2'$.
\end{proof}

\subsection{Small-shift Type I}\label{subsec:small-typeI}
\begin{proposition}[Small-shift Type I]\label{prop:small-I}
LWW Section~3 with the following divisor and parity substitutions
gives $T_1(F_h),T_2(F_h)\ll_A xL^{-A}$ uniformly in $h$.
\end{proposition}
\begin{proof}
The sifted $q$ satisfies $(q,h)=1$.
After expanding $R_Z$, let $g=(r,h)$ with $r$ odd squarefree.
The list divisor $d$ is disjoint from $rh$.
Apply \eqref{eq:mu-dg} and write
\begin{equation}\label{eq:small-I-substitution}
 qrj+h=dgat,\qquad
 t\equiv(h/g)(da)^{-1}\pmod{qr/g},\qquad t\ge1.
\end{equation}
The inverse exists because $(da,qr/g)=1$.
The class is reduced since $(h/g,r/g)=(h,q)=1$ and
$(g,qr/g)=1$. The primes of $a$ may lie in $g$, so no term is
removed by a condition $(g,da)=1$.
The M\"obius function is evaluated at $t$ itself in this substitution;
factoring out $g$ after evaluating it would instead require
$\mu(gt')=\mu(g)\mu(t')\ind_{(t',g)=1}$.

For even $h$, $j$ odd is equivalent to $t$ odd, and the class
modulo $2qr/g$ is reduced.
For odd $h$, $j$ odd is equivalent to $t$ even. Put $t=2t''$;
then
\begin{equation}\label{eq:small-I-parity}
 t''\equiv(h/g)(2da)^{-1}\pmod{qr/g},\qquad
 \mu(2t'')=-\mu(t'')\ind_{2\nmid t''}.
\end{equation}
The last indicator makes the class modulo $2$ reduced, so in both
cases the M\"obius sum is over a reduced class modulo $2qr/g$.
Apply (3.1) with the weighted grouping of LWW p.~13 separately
at each $g$-dependent level $x/(dg)$ (and $x/(dga)$ in the small-$a$
part). Then sum $\sum_{g\mid\rad(h)}1/g\ll\log L$.
The extra Rankin factor
$\prod_{p\mid g}(1-p^{-1/20})^{-1}\le L^{O(1)}$ is absorbed by
$x^{-1/50}$ as in \eqref{eq:I-Rankin}.
The positive interval for $p+h\in(x+h,2x+h]$ is handled by the
maximum over endpoints in the same estimate.
\end{proof}

\begin{remark}[Other Type I parametrizations]\label{rem:other-I}
For comparison, factoring $t=gt'$ after the original substitution
requires the factor $\ind_{(t',g)=1}$ in
$\mu(gt')=\mu(g)\mu(t')\ind_{(t',g)=1}$.
One may split into at most $\ph(g)\le|h|$ reduced classes modulo
$qr$ and absorb $L^C$ by increasing the distribution precision.
Another exact expansion is
\[
 R_Z(k)=\ind_{(k,h)=1}R'_Z(k),\qquad
 R'_Z(k)=\ind_{\text{each prime divisor of $k$ exceeds $Z$ or divides $h$}}.
\]
Expanding only $R'_Z$ leaves $r$ coprime to $h$; the remaining
condition on $j$ can be expanded over $e\mid\rad(h)$, with
$2^{\omega_{\rm all}(h)}\le L^{O(1)}$ terms. The proof above uses
\eqref{eq:small-I-substitution} instead.
\end{remark}

\subsection{Removal of weights and assembly}\label{subsec:small-assembly}
\begin{proof}[Proof of Theorem~\ref{thm:H}]
For weight removal, $d\mid P_0$ and $(h,d)=1$ give
\begin{equation}\label{eq:small-weight-CRT}
 \frac1{\ph(P_0)}\sum_{\substack{a\bmod P_0\\(a,P_0)=1}}
          \ind_{d\mid a+h}=\frac1{\ph(d)}.
\end{equation}
Bombieri--Vinogradov uses the reduced class $p\equiv-h\pmod d$.
Use the same mean normalization as in Proposition~\ref{prop:weights}.

For squarefree density use $1\le p+h\le x+|h|$.
If $p\nmid h$ and $d^2\mid p+h$, then $(d,h)=1$; otherwise a
common prime divisor would equal $p\mid h$, which accounts for
at most $\omega_{\rm all}(h)$ primes.
Siegel--Walfisz applies uniformly to the reduced class
$p\equiv-h\pmod{d^2}$ for $d\le L^{A+2}$.
For the tail, restricting to $d\le\sqrt{x+|h|}$ gives
\begin{equation}\label{eq:small-squarefree-tail}
 L\sum_{d>L^{A+2}}\frac{x+|h|}{d^2}
       +O(\sqrt{x+|h|}\,L),
\end{equation}
and the extra endpoint term is absorbed at any fixed precision.
Thus
\begin{equation}\label{eq:small-squarefree}
 \#\{p\le x:p+h\ge1,\ \mu^2(p+h)=1\}
       =(A_h+o(1))\pi(x),\qquad A_h\ge0.3739,
\end{equation}
uniformly in $h$.

Assemble Propositions~\ref{prop:small-II}, \ref{prop:small-major}
and~\ref{prop:small-I}, using
$|T_Y|^2\ll(X/Y)Q_Y$ and $Q_Y\ll XYL^{-D^*}$.
The dyadic summation and normalized weight removal give the weighted
cancellation, and the bounded-target comparison
\eqref{eq:unweight}, with $\GN$ replaced by $x$, gives its
unweighted form. The negligible lower prime range includes all
$p+h\le0$, at most $O(L^C)$ primes; the domain in
\eqref{eq:H} excludes them throughout.
When summing dyadic intervals with lower endpoint $\gg x/L^B$,
their logarithms are $L+O_B(\log L)$, so the fixed-$C$ range and
all bounds above remain uniform. Apply the sign identity
\eqref{eq:mu-sign-identity} on positive arguments to obtain
\eqref{eq:H}. This uses LWW Sections~2.3, 3, 5 and 6, as well as
$\mathrm H2'$, exactly as in the statement.
\end{proof}

\section{Liouville versions}\label{sec:liouville}
The function $\lambda(n)=(-1)^{\Omega(n)}$ is completely
multiplicative and takes values in $\{-1,1\}$ for every $n\ge1$.
We retain the Vaughan coefficient $\mu(m)R_Z(m)$, since it comes
from decomposing the prime variable.
For the shifted target define explicitly
\begin{equation}\label{eq:F-lambda}
 \begin{split}
 F_{\lambda,G}(n)&=\begin{cases}
 \lambda(\GN-n)W(\GN-n),&1\le n<\GN,\\0,&\text{otherwise},\end{cases}\\
 F_{\lambda,h}(n)&=\begin{cases}
 \lambda(n+h)W(n+h),&n\ge1,\ n+h\ge1,\\0,&\text{otherwise}.
 \end{cases}
 \end{split}
\end{equation}

For $\kappa=\pm1$, the three statements of
Theorem~\ref{thm:lambda} are explicitly
\begin{align}
 \sum_{p<\GN}\lambda(\GN-p)&=o(\pi(\GN)),&
 \#\{p<\GN:\lambda(\GN-p)=\kappa\}
       &=\left(\tfrac12+o(1)\right)\pi(\GN),\label{eq:Glambda}\\
 \sum_{\substack{p\le x\\p+h\ge1}}\lambda(p+h)&=o(\pi(x)),&
 \#\{p\le x:p+h\ge1,\ \lambda(p+h)=\kappa\}
       &=\left(\tfrac12+o(1)\right)\pi(x).\label{eq:Hlambda}
\end{align}
The former holds for all even $\GN$ under the hypotheses of
Theorem~\ref{thm:G}, and for the fixed-$C$ small-radical class
under those of Theorem~\ref{thm:S}; the latter is uniform for
$0<|h|\le(\log x)^C$ under those of Theorem~\ref{thm:H}.
LWW Section~1, pp.~1--3, mentions the Liouville counterpart; we
supply the expansion and Type~I substitutions needed for it.

\subsection{Expansion of the weight}\label{subsec:lambda-expansion}
The exact identities are
\begin{equation}\label{eq:lambda-weight-expansion}
 \begin{split}
 \omega_i(n)\rho^{\omega_i(n)-1}
    &=\sum_{\substack{p\in\mathcal P_i\\p\mid n}}\rho^{\omega_i(n)-1},\\
 W(n)&=\frac1{\prod_iV_i}\sum_{\substack{a\in\Aset\\a\mid n}}
                                \rho^{\omega(n)-K},\\
 \omega(ah)-K&=\omega(h)-s(a,h),\qquad
 s(a,h)=\#\{i:p_i\mid h\}\in\{0,\ldots,K\},
 \end{split}
\end{equation}
where $\omega_i(ah)=\omega_i(h)+\ind_{p_i\nmid h}$ and
$s(a,h)=0$ precisely when $(a,h)=1$.

\begin{lemma}[Liouville expansion]\label{lem:lambda-expansion}
Let $w$ denote $mn+2$, $\GN-mn$ or $mn+h$, with the positive
argument convention used above and the corresponding coefficients.
Then
\begin{equation}\label{eq:lambda-expansion}
 \begin{split}
 \sum_{\substack{m,n\\w\ge1}}\alpha_m\beta_n\lambda(w)W(w)
 ={}&\frac{(-1)^K}{\prod_iV_i}\sum_{\ell\ge1}\lambda(\ell)
                  \rho^{\omega(\ell)}
 \sum_{\substack{a\in\Aset,m,n\\w=a\ell,(a,\ell)=1}}\alpha_m\beta_n\\
 &+O_{K,\rho}(XL e^{-L^{0.1}}).
 \end{split}
\end{equation}
For a fixed positive $w$, the discrepancy before multiplication by
$\alpha_m\beta_n$ is exactly
\begin{equation}\label{eq:lambda-discrepancy}
 \frac{\lambda(w)}{\prod_iV_i}
 \sum_{\substack{a\in\Aset,a\mid w\\s(a,w/a)\ge1}}
                              \rho^{\omega(w)-K}.
\end{equation}
It lies in $\lambda(w)[0,W(w)]$; every summand weight is at most $1$.
\end{lemma}
\begin{proof}
For $a\in\Aset$, $\lambda(a)=(-1)^K$, and complete
multiplicativity gives $\lambda(w)=(-1)^K\lambda(w/a)$ for every
$a\mid w$, with no coprimality assumption.
On $(a,w/a)=1$, \eqref{eq:lambda-weight-expansion} gives exactly
the main term. The omitted summands are precisely
\eqref{eq:lambda-discrepancy}.
They have some $p_i^2\mid w$, and $\omega(w)\ge K$ makes their
weights at most $1$. Their absolute value is bounded, with constant
$1$, by
\[
 \frac1{\prod_iV_i}\sum_{\substack{a\in\Aset,a\mid w\\(a,w/a)>1}}1.
\]
Since $W(w)\ll_{K,\rho}1$, the absolute discrepancy is at most
\[
 C_{K,\rho}\sum_{p\in\cup_i\mathcal P_i}\ind_{p^2\mid w}.
\]
For fixed $m$ on the coefficient support, $(m,p)=1$ and the
condition $p^2\mid w$ selects one class of $n$ modulo $p^2$.
Since $N\gg p^2$, its count is at most $N/p^2+1\ll N/p^2$.
Summing with $|\alpha_m|\le1$, $|\beta_n|\ll L$ and
$\sum_{p\in\cup_i\mathcal P_i}p^{-2}\ll e^{-L^{0.1}}$
gives the claimed error for all three positive shifted arguments.
This is the bounded-weight square-divisibility argument of
LWW (4.3), p.~15.

If the restriction $(a,\ell)=1$ is then removed before Cauchy,
the added terms have $s\ge1$ and
$\rho^{\omega(\ell)}\le\rho^{\omega(w)-K}$, by
\eqref{eq:lambda-weight-expansion}. They are bounded by the same
count, at an additional fixed constant.
Cauchy uses only $|\lambda(\ell)\rho^{\omega(\ell)}|\le1$.
For the Goldbach shift use $Q_Y^+\le\QQ$ in
\eqref{eq:cauchy-all}; for small shifts use \eqref{eq:small-cauchy}.
In each case the resulting square contains no target function.
\end{proof}

\subsection{Bombieri--Vinogradov for Liouville}\label{subsec:lambda-BV}
The M\"obius input is LWW (3.1), p.~10: for every $A>0$ there
exists $B=B(A)$ such that
\begin{equation}\label{eq:BVmu}
 \sum_{q\le x^{1/2}L^{-B}}
 \max_{(c,q)=1}\max_{y\le x}
 \left|\sum_{\substack{n\le y\\n\equiv c\ (q)}}\mu(n)\right|
 \ll_A xL^{-A}.
\end{equation}
\begin{lemma}[Liouville distribution]\label{lem:BVlambda}
For every $A>0$ there is $B'(A)$ such that
\begin{equation}\label{eq:BVlambda}
 \sum_{q\le x^{1/2}L^{-B'(A)}}
 \max_{(c,q)=1}\max_{y\le x}
 \left|\sum_{\substack{n\le y\\n\equiv c\ (q)}}\lambda(n)\right|
 \ll_A xL^{-A}.
\end{equation}
One may take $B'(A)=\max(B(A),A)+A+1$.
\end{lemma}
\begin{proof}
The multiplicative identity
\begin{equation}\label{eq:lambda-mu-convolution}
 \lambda(n)=\sum_{d^2\mid n}\mu(n/d^2)
\end{equation}
holds prime by prime: at $p^k$ the sum is $1$ for even $k$ and
$-1$ for odd $k$.
For $(c,q)=1$,
\[
 \sum_{\substack{n\le y\\n\equiv c\ (q)}}\lambda(n)
 =\sum_{d\le\sqrt y}\sum_{\substack{m\le y/d^2\\d^2m\equiv c\ (q)}}\mu(m).
\]
There is no solution when $(d,q)>1$; otherwise the class is
$m\equiv c\bar d^{\,2}\pmod q$, which is reduced.
Put $D_0=L^{A+1}$ and $Q=x^{1/2}L^{-B'}$.
For $d\le D_0$, use the discrepancy
\[
 E_\mu(x/d^2;q)=\max_{(c',q)=1}\max_{z\le x/d^2}
 \left|\sum_{\substack{m\le z\\m\equiv c'\ (q)}}\mu(m)\right|.
\]
For $d>D_0$ use $y/(d^2q)+1$.
The tail summed over $q\le Q$ is
\begin{equation}\label{eq:lambda-BV-tail}
 \ll xL/D_0+Q\sqrt x\le xL^{-A}+xL^{-B'}.
\end{equation}
For the head, $x_d=x/d^2\ge xL^{-2A-2}$ and
$\log x_d\ge L/2$. Condition $B'\ge B(A)+A+1$ ensures
$Q\le\sqrt{x_d}(\log x_d)^{-B(A)}$.
Apply \eqref{eq:BVmu} at $x_d$ and sum
$\sum_{d\le D_0}xd^{-2}L^{-A}\ll_A xL^{-A}$.
The stated $B'(A)$ also absorbs the tail.
\end{proof}

\begin{lemma}[Weighted Liouville distribution]\label{lem:BVlambda-weighted}
Put
\[
 E_\lambda(Y;m)=\max_{(c,m)=1}\max_{z\le Y}
 \left|\sum_{\substack{t\le z\\t\equiv c\ (m)}}\lambda(t)\right|.
\]
For fixed $A_{\rm BV}>2A+4$ and
$Q\le Y^{1/2}L^{-B'(A_{\rm BV})}$,
\begin{equation}\label{eq:BVlambda-weighted}
 \sum_{m\le Q}\tau(m)E_\lambda(Y;m)\ll_A YL^{-A}.
\end{equation}
Here $\log Y\asymp L$ in the applications, with fixed comparison
constants and a corresponding fixed precision margin.
\end{lemma}
\begin{proof}
The trivial bound is $E_\lambda(Y;m)\ll Y/m+1$.
As in the Cauchy step of LWW Section~3, p.~13,
\begin{equation}\label{eq:BVlambda-Cauchy}
 \left(\sum\tau E\right)^2
 \le\left(\sum E\right)\left(\sum\tau^2E\right),\quad
 \sum E\ll YL^{-A_{\rm BV}},\quad
 \sum\tau^2E\ll YL^4+QL^3.
\end{equation}
Use Lemma~\ref{lem:BVlambda} at precision $A_{\rm BV}$, rather
than $A$, and take square roots.
\end{proof}

\subsection{Type I for Liouville}\label{subsec:lambda-typeI}
\begin{proposition}[Liouville Type I]\label{prop:lambda-I}
With the corresponding positive-argument functions $F_\lambda$,
$T_1(F_\lambda),T_2(F_\lambda)\ll_{A,K,\rho}xL^{-A}$ for
shift $+2$, Goldbach shift $\GN$, and small shift $h$ under the
same uniformity ranges as above.
\end{proposition}
\begin{proof}
The truncated inclusion--exclusion for $R_Z$, divisor expansion of
$W$ and truncation error (3.3) in LWW Section~3, pp.~10--12,
use only $|\alpha_q|\le L$, boundedness of $F$, and counts in
residue classes. For shift $+2$ the remaining expression is
\begin{equation}\label{eq:lambda-I-shape}
 \sum_q\alpha_q\sum_r\mu(r)
 \sum_{\substack{d\\(d,r)=1}}c_d
 \sum_{\substack{x<qk\le y\\2\nmid k, r\mid k\\d\mid qk+2}}
                                      \lambda(qk+2),
 \qquad y\in[x,2x].
\end{equation}
Write $k=rj$. Now $d$ odd and $d\mid qrj+2$ imply $(d,qr)=1$.
Replace (3.4) by complete multiplicativity:
\begin{equation}\label{eq:lambda-I-two}
 \lambda(dv)=\lambda(d)\lambda(v),\qquad
 qrj+2=dt,\qquad
 \lambda(d)\sum_{\substack{(x+2)/d<t\le(y+2)/d\\
                         dt\equiv qr+2\ (2qr)}}\lambda(t).
\end{equation}
The class is reduced since $(d,2qr)=(qr+2,2qr)=1$.
There is no auxiliary $a$-sum. The divisor satisfies
$d\le\exp(O_K(R_I L^{0.2}))=x^{o(1)}$, where
$R_I=\lceil C'\log L\rceil$ is the Type~I truncation depth.
For large $x$ this gives $d\le x^{1/20}$; for the Goldbach
substitution below, also $dg\le x^{1/20}$ since $g\le r=x^{o(1)}$.
Thus the small-divisor case $da\le x^{1/20}$ with $a=1$ covers the whole sum,
and no Rankin branch is needed.

Group by $qr$ with coefficient at most $L\tau(qr)$.
For $Y=(2x+2)/d\ge2x^{19/20}$ and
$2qr\le x^{2/5+o(1)}\le Y^{1/2}L^{-B'(A_{\rm BV})}$,
Lemma~\ref{lem:BVlambda-weighted} gives
\begin{equation}\label{eq:lambda-I-bound}
 \sum_d|c_d|\sum_{q,r}|\alpha_q|E_\lambda((2x+2)/d;2qr)
 \ll xL^{-A}\sum_d|c_d|/d\ll_{K,\rho}xL^{-A},
\end{equation}
using $\sum_d|c_d|/d\ll_{K,\rho}1$ from LWW Section~3, p.~11.
The $T_1$ partial summation of p.~14 is unchanged.

For the Goldbach shift, let $g=(r,\GN)$ and write
\begin{equation}\label{eq:lambda-I-G}
 \GN-qrj=dgt,\qquad
 \lambda(\GN-qrj)=\lambda(dg)\lambda(t),\qquad
 t\equiv(\GN/g)d^{-1}\pmod{qr/g}.
\end{equation}
The class is reduced by the coprimality checks in
Proposition~\ref{prop:I}, and $j$ odd is equivalent to $t$ odd.
Use modulus $2qr/g$ and level $Y=\GN/(dg)\ge x^{19/20}$;
the modulus is $\le x^{2/5+o(1)}$.
Carry out the grouping separately for each $g$ because this level
depends on $g$, then sum $\sum_{g\mid\rad(\GN)}1/g\ll\log L$.
The endpoints are those of Section~\ref{sec:typeI} with $a=1$.

For the shift $h$, put $g=(r,h)$ and use
\begin{equation}\label{eq:lambda-I-h}
 qrj+h=dgt,\qquad \lambda(qrj+h)=\lambda(dg)\lambda(t),\qquad
 t\equiv(h/g)d^{-1}\pmod{qr/g}.
\end{equation}
For even $h$ the odd $t$ class modulo $2qr/g$ is reduced.
For odd $h$, put $t=2t''$:
\begin{equation}\label{eq:lambda-I-odd}
 \lambda(2t'')=-\lambda(t''),\qquad
 t''\equiv(h/g)(2d)^{-1}\pmod{qr/g}.
\end{equation}
There is no parity condition on $t''$ here.
Apply Lemma~\ref{lem:BVlambda-weighted} at modulus $qr/g$ and
$Y=(2x+h)/(2dg)$; handle each $g$ separately and then sum $1/g$.
All cofactors are positive and $0$ is excluded.
\end{proof}

\subsection{Uses of the target function}\label{subsec:lambda-inventory}
\begin{center}
\small
\begin{longtable}{@{}>{\raggedright\arraybackslash}p{.07\textwidth}>{\raggedright\arraybackslash}p{.39\textwidth}>{\raggedright\arraybackslash}p{.46\textwidth}@{}}
\caption{Target-function substitutions.}\label{tab:lambda-inventory}\\
\toprule
Item&LWW input or operation&Liouville substitution\\\midrule
\endfirsthead
\toprule Item&LWW input or operation&Liouville substitution\\\midrule
\endhead
1&(4.3), p.~15; weight expansion&Lemma~\ref{lem:lambda-expansion}, including the square-overlap count.\\
2&(4.4)--(4.5), p.~15; Cauchy&$|\lambda(\ell)\rho^{\omega(\ell)}|\le1$; the Goldbach square is \eqref{eq:cauchy-all}.\\
3&Sections~4.2--4.7, pp.~17--39; Appendix~A, pp.~47--51&The square has no target function; the same moment and root arguments apply.\\
4&Section~5, pp.~39--45&The M\"obius coefficient in $M$ remains $\mu(m)R_Z(m)$ from Vaughan's identity.\\
5&(3.1), p.~10&Lemma~\ref{lem:BVlambda}.\\
6&(3.4), p.~12, and (3.5), p.~13&$\lambda(dv)=\lambda(d)\lambda(v)$; no $a$-sum or Rankin branch.\\
7&Weighted grouping, p.~13&Lemma~\ref{lem:BVlambda-weighted}, using the higher fixed precision.\\
8--10&(3.3) truncations; Proposition~2.1; Proposition~2.3 and weight removal&Only boundedness of the target and residue counts; use the same positive means $m_j$.\\
11&Proposition~6.1, pp.~46--47&No squarefree-density input is required.\\
12&Final sign assembly&Use \eqref{eq:lambda-sign-identity}.\\
\bottomrule
\end{longtable}
\end{center}
Other M\"obius and squarefree factors in LWW occur in the divisor
combinatorics of Proposition~2.3 (pp.~7--9), Bombieri--Vinogradov
for primes (p.~9), squarefree $r,c_d$ in the Type~I expansion
(pp.~11--14), the relative-coordinate/list structure (pp.~22--26,
35), the Vaughan coefficient, and the appendix sieve (p.~50).
They do not refer to the target sign and are retained.
The additional unit discards in Sections~\ref{sec:setting}
and~\ref{sec:typeI} use only boundedness.

\subsection{Sign densities and assembly}\label{subsec:lambda-assembly}
\begin{proof}[Proof of Theorem~\ref{thm:lambda}]
For any finite prime set $\mathcal P$ and positive $n(p)$,
\begin{equation}\label{eq:lambda-sign-identity}
 \#\{p\in\mathcal P:\lambda(n(p))=\kappa\}
 =\tfrac12\left(|\mathcal P|+
          \kappa\sum_{p\in\mathcal P}\lambda(n(p))\right),\quad\kappa=\pm1.
\end{equation}
Use $\mathcal P=\{p<\GN\}$, of size $\pi(\GN)$ for even
$\GN>2$, or $\mathcal P=\{p\le x:p+h\ge1\}$, of size
$\pi(x)+O(L^C)$.

In the Goldbach assembly, use Lemma~\ref{lem:lambda-expansion} and
Proposition~\ref{prop:lambda-I} for the expansion and Type~I sums.
The block estimate, major arcs and Theorem~\ref{thm:A} are independent
of the target. The same normalized weight removal gives the unweighted
cancellation. For small shifts use the identical substitutions in
Section~\ref{subsec:small-assembly}.
Lemma~\ref{lem:BVlambda} changes the fixed distribution exponent
$B(A)$ to $B'(A)$; all choices still precede the large parameter.
Uniformity in $\GN$ or $|h|\le L^C$ and the respective hypotheses
are unchanged. Identity~\eqref{eq:lambda-sign-identity} proves the
sign densities without Proposition~\ref{prop:squarefree}.

On squarefree values $\lambda=\mu$.
Thus the nonsquarefree values in the Goldbach setting have density
$(1-A_{\GN})/2+o(1)$ for each Liouville sign, consistent with
\eqref{eq:G} and \eqref{eq:Glambda}.
The errors remain ineffective through the cancellation argument.
\end{proof}

\appendix
\section{Changes relative to LWW, Sections 4.1--4.7 and Appendix A}
\label{app:dictionary}
All LWW references in Table~\ref{tab:dictionary} are to version~1,
14 October 2026. Page numbers are the numbered pages of that paper.
Its arguments are used as written with these replacements.
The square being represented is the full all-integer nonnegative
square \eqref{eq:cauchy-all}, whose expansion is
\begin{equation}\label{eq:all-expanded}
 \QQ=\prod_iV_i^{-2}
 \sum_{\substack{a,b,m,n,r,s\\
    (\GN-mn)/a=(\GN-rs)/b\in\Z\setminus\{0\}}}
 \eta(a/Y)\eta(b/Y)\alpha_m\overline{\alpha_r}
                    \beta_n\overline{\beta_s}.
\end{equation}
After the common-list removal, \eqref{eq:determinant} is exactly
its determinant condition. Endpoint pairing uses
\eqref{eq:endpoint-conjugation} with a conjugate-linear initial
entry. The moment normalization remains $d_0YMN$, and Cauchy uses
$\GN/Y$ separately.

\small
\begin{longtable}{@{}>{\raggedright\arraybackslash}p{.045\textwidth}>{\raggedright\arraybackslash}p{.235\textwidth}>{\raggedright\arraybackslash}p{.655\textwidth}@{}}
\caption{The block dictionary.}\label{tab:dictionary}\\
\toprule
Row&LWW location&Replacement or retained input\\\midrule
\endfirsthead
\toprule Row&LWW location&Replacement or retained input\\\midrule
\endhead
1\phantomsection\label{dict:01}&Section~4.1, p.~14; (4.3)--(4.5), p.~15&
Separate the product cutoff as there, with $F=0$ off positive shifted arguments.
Use the nonnegative enlargement \eqref{eq:cauchy-all} and its exact
integer-cofactor expansion \eqref{eq:all-expanded}. Its common range
is $|h|\ll\GN/Y$, and $|T_Y|^2\ll(\GN/Y)\QQ$.\\
2\phantomsection\label{dict:02}&(4.8)--(4.9), p.~16; (4.24)--(4.25), p.~21&
Set $t'=t/\GN$, use $\psi_1(t/(\GN Y))=\psi_1(t'/Y)$ and
$\MN$. The $\GN$ branches cancel the Jacobian $1/\GN$ in
Lemma~\ref{lem:cover}.\\
3\phantomsection\label{dict:03}&Proposition~4.3, pp.~20--21; Section~4.6, p.~35&
Use labelled index-$\GN$ blocks and $P_j=\gN(P)z_j$ with
$\gN=gE_{\GN}U_0$.
At endpoints use $f(Db m,s)=\overline{\alpha_m}\beta_s$, so
$\overline{f(P)}f(Q)=\alpha_m\overline{\alpha_r}\beta_n\overline{\beta_s}$.\\
4\phantomsection\label{dict:04}&Continuous extension, p.~22&
Keep the entry $\GN/u$ exactly in \eqref{eq:continuous};
$\tauN(z)=z_1+\beta z_2$ replaces $z_1-2rz_2$.\\
5\phantomsection\label{dict:05}&Path bound, p.~23&
Lemma~\ref{lem:coordinates} directly gives
$|z_2|\le64d_0YMN/\GN$, $|z_1|\le16+|z_2|$.
It implies the needed $|z|\ll Rd_0Y$ and coefficient mass
$\exp(O_K(L^{0.73}))$.\\
6\phantomsection\label{dict:06}&(4.30), (4.32), pp.~23--24&
The second boundary is
$E_2=t_0\tauN+z_2\GN/(d_0YMNs)$.
The extra term is $O(1)$, as is its $s$ derivative;
$\partial_{t_0}E_2=\tauN\in[1/16,16]$ and
$\partial_\beta E_2=t_0z_2=\exp(O(L^{0.21}))$.
The neighborhood volume is $x^{-c\delta/8+o(1)}$.\\
7\phantomsection\label{dict:07}&(4.31), p.~24&
$\partial_\beta(\tauN(w)/\tauN(z))=1/\tauN(z)^2$ for
$\det(z,w)=1$, replacing $-2/\tau(z)^2$.\\
8\phantomsection\label{dict:08}&Boundary cubes, p.~24&
Use $(s,t_0,\beta)$, the same transverse derivatives and
$x^{o(1)}$ choices. Sum the interval-uniform errors of Lemma~\ref{lem:R}
over cubes before the moment limit.\\
9\phantomsection\label{dict:09}&(4.33)--(4.34), pp.~24--25&
Use the primitive relative box of radius $2C_{\rm pos}Rd_0Y$
and the same $\tau$ window. The multiplier is
$\chi_{\GN}=\ind_{\gN z\in\Omega}$, of norm at most $1$.
Endpoint unit support replaces LWW's $z_1$-odd test;
that test is not implied for the normalized coordinates used here.\\
10\phantomsection\label{dict:10}&(4.40)--(4.41), pp.~27--28&
$z'=j_1w_z+hz$ and
$\tauN(z')/\tauN(z)=h+j_1\tauN(w_z)/\tauN(z)$.
The allowed ratio strip leaves $O(1)$ integers $h$ per $j_1$.\\
11\phantomsection\label{dict:11}&Slope-cell identity, p.~28&
$z'_2/\tauN(z')-z_2/\tauN(z)=\det(z,z')/(\tauN(z)\tauN(z'))$.
An edge has reduced determinant $O(d_0Y)$; cells of width $d_0Y$
have bounded interaction overlap. A Goldbach block meets $O(1)$ such cells.\\
12\phantomsection\label{dict:12}&(4.42)--(4.46), pp.~28--30&
Use the shared-prime lattice, Minkowski box $A'B'\asymp Y$ and
the same signed kernel
$Y^{-1}\psi_1(\det(k,\ell)/Y)\int_{\T\setminus\Mone}e(\theta\det(k,\ell))\,d\theta$.
Its $\alpha=D\tauN(w_*)/\tauN(z_*)$ is $D\beta_{Q_*}$ modulo $1$.
The enclosing region has $|\ell|\ll Y$ and bounded $\tau$ ratios.\\
13\phantomsection\label{dict:13}&(4.45), p.~29&
$\chi(B_Dk)$ is a bounded endpoint multiplier. Sharp or nonrectangular
position restrictions preserve its contraction property.\\
14\phantomsection\label{dict:14}&(4.47)--(4.48), p.~30&
$j_1=P_{\rm full}-Dp_*ZT_f=\det(z,z')=D(b-a)$;
$\tauN(z')/\tauN(z)\in[1/16,16]$.
The second good-state test is applied to the intrinsic ratio at
the source $g_{\GN}(P)z$.\\
15\phantomsection\label{dict:15}&(4.49)--(4.50), p.~31&
The reduced integer box has the same nonzero linear congruence modulo
a list prime. List primes are coprime to $\GN$ and to the chosen shared $D$.\\
16\phantomsection\label{dict:16}&Endpoint primitivity, p.~35&
A common prime of $m,s$ is $>Z$ and coprime to $\GN$.
It divides $b-a$, while $0<|b-a|<4Y<Z$ by the retained cross-edge
exclusion. The cofactor sign is irrelevant.\\
17\phantomsection\label{dict:17}&Section~4.6 partition, p.~36&
Use labelled blocks with $O(d_0Y)$ lattice-primitive positions in the
application of Lemma~\ref{lem:block-size}, with no interaction between
blocks. The disjoint endpoint union is restricted to units as in
\eqref{eq:disjoint-union}.\\
18\phantomsection\label{dict:18}&(4.63)--(4.68), pp.~36--37&
Use the positive principal submatrix on
$\CC_{\GN}\cap\Lambda_\varrho$ in the enlarged ambient space.
Spectral H\"older precedes root restriction;
$\lambda_{\max}\le d_\varrho$ and $(d_0Y)^{1/(2R)}=O(1)$.
Only the restricted-root trace is summed.\\
19\phantomsection\label{dict:19}&(4.70), p.~37&
Exact-edge $O(1)$ and comparison-edge $O(Y)$ counts are in $t'$.
The bad-state ratio measure is unchanged.\\
20\phantomsection\label{dict:20}&(4.74)--(4.76), p.~38; (4.6), p.~16&
For fixed $k_0=mn$, $\ell_0=rs$ is in the class
$ba^{-1}k_0\bmod\GN$ in an interval of length $O(\GN)$, giving
$O(1)$ choices. Raw absolute divisor-square counts use the translated
cofactor interval or $|h|\ll\GN/Y$, at a cost $\GN/X\ll L^B$.\\
21\phantomsection\label{dict:21}&Appendix~A, swapped case, p.~50&
Retain $\GN/(uv)$, with $\GN/(uv)\ll L^B/(d_0Y)$.
Use both $k=\lfloor\beta+\GN/(uv)\rfloor\in\{0,1\}$ branches,
the separate matrix residues and the cell envelopes of
\eqref{eq:swapped}--\eqref{eq:swapped-residue}.\\
22\phantomsection\label{dict:22}&Lemma~A.1 and (A.5), pp.~47--51&
Use Lemma~\ref{lem:R}, the determinant-$\GN$ matrix coset,
$c_{\GN}=\zeta(2)^{-1}\prod_{p\mid\GN}(p-1)/(p+1)$ and
line masses $1/(p+1)$. Its factor at $2$ is already included.
Summing matrix errors costs $S^3=x^{o(1)}$ and decreases the final
fixed line-error exponent.\\
23\phantomsection\label{dict:23}&Arc measure, p.~16&
Haar measure under $\theta\mapsto\GN\theta$ gives
$|\MN|=|\Mone|\ll L^{3A_0}/Y$.\\
\bottomrule
\end{longtable}
\normalsize

\bibliographystyle{plainnat}
\bibliography{references}
\section*{Use of AI}
The proofs and the text of this paper were produced with AI systems
(GPT-6.1 Sol, Claude Fable~5.1 and Claude Opus~5.5) under the author's
direction. The proofs were checked by independent reviews with AI
models, and the algebraic identities on finite instances by the
programs in the ancillary files. The author takes responsibility for the content.
\end{document}
